\documentclass[a4paper,11pt,oneside]{scrbook}

% Pengaturan Font dan Bahasa untuk LuaLaTeX
\usepackage{fontspec}
\usepackage[indonesian]{babel}
\usepackage{amsmath,amssymb,amsfonts}
\usepackage{graphicx}
\usepackage[space]{grffile} % Mendukung nama file gambar yang mengandung spasi seperti "Tugas 3-001.png"
\usepackage{adjustbox}      % Mencegah rumus panjang keluar dari margin kanan
\usepackage{tikz}           % Mendukung ilustrasi vektor jika diperlukan
\usepackage{xcolor}
\usepackage{listings}
\usepackage{fancyvrb}
\usepackage{enumitem}
\usepackage{geometry}
\usepackage{hyperref}

\geometry{
    a4paper,
    top=2.5cm,
    bottom=2.5cm,
    left=2.5cm,
    right=2.5cm
}

\setlist[itemize]{label=-, topsep=4pt, itemsep=2pt}

% Lokasi folder gambar
\graphicspath{{./}{images/}}

% Definisi warna ala Euler Math Toolbox (EMT)
\definecolor{eulerpromptcolor}{rgb}{0.6,0.0,0.0}
\definecolor{euleroutputcolor}{rgb}{0.15,0.15,0.45}
\definecolor{eulerbg}{rgb}{0.97,0.97,0.98}
\definecolor{eulerborder}{rgb}{0.82,0.85,0.90}

\hypersetup{
    colorlinks=true,
    linkcolor=euleroutputcolor,
    urlcolor=eulerpromptcolor,
    pdftitle={Menggambar Grafik 2D dengan EMT}
}

% Pengaturan Listings untuk Kode EMT (rata kiri tanpa indentasi berlebih)
\lstset{
    basicstyle=\ttfamily\small,
    breaklines=true,
    breakatwhitespace=false,
    breakindent=0pt,
    breakautoindent=false,
    columns=fixed,
    basewidth=0.52em,
    keepspaces=true,
    showstringspaces=false,
    frame=leftline,
    framerule=1.5pt,
    framexleftmargin=6pt,
    xleftmargin=8pt,
    xrightmargin=0pt,
    resetmargins=true,
    rulecolor=\color{eulerborder},
    backgroundcolor=\color{eulerbg},
    aboveskip=6pt,
    belowskip=4pt
}

% Lingkungan khusus pengganti paket eumat agar kompatibel penuh dengan LuaLaTeX
\newenvironment{eulernotebook}{}{}

\newcommand{\eulerheading}[1]{\chapter{#1}}
\newcommand{\eulersubheading}[1]{\section{#1}}

\newenvironment{eulercomment}{\par\medskip\noindent}{\par\medskip}

% Otomatis menyesuaikan rumus yang terlalu lebar agar tidak terpotong margin kanan
\newenvironment{eulerformula}{%
    \par\medskip\noindent
    \renewcommand{\[}{\begin{center}\begin{adjustbox}{max width=\textwidth}$\displaystyle}%
    \renewcommand{\]}{$\end{adjustbox}\end{center}}%
}{\par\medskip}

\lstnewenvironment{eulerprompt}{
    \lstset{basicstyle=\ttfamily\small\color{eulerpromptcolor}}
}{}

\lstnewenvironment{euleroutput}{
    \lstset{
        basicstyle=\ttfamily\footnotesize\color{euleroutputcolor},
        backgroundcolor=\color{white},
        rulecolor=\color{lightgray}
    }
}{}

\lstnewenvironment{eulerudf}{
    \lstset{basicstyle=\ttfamily\small\color{teal!70!black}}
}{}

\lstnewenvironment{eulerttcomment}{
    \lstset{basicstyle=\ttfamily\small\color{darkgray}}
}{}

% Perintah penampil gambar yang menjaga proporsi rumus kecil (#1 < 10) maupun plot grafik (#1 >= 10)
\newcommand{\eulerimg}[2]{%
    \ifnum#1>0
        \begin{center}
            \IfFileExists{#2}{%
                \ifnum#1<10
                    \includegraphics[width=0.88\textwidth,height=#1\dimexpr 1.75ex\relax,keepaspectratio]{#2}%
                \else
                    \includegraphics[width=0.72\textwidth,height=#1\dimexpr 1.85ex\relax,keepaspectratio]{#2}%
                \fi
            }{%
                \fbox{\parbox{0.7\textwidth}{\centering\small\texttt{Gambar tidak ditemukan: \detokenize{#2}}\\ \footnotesize Pastikan folder \texttt{images/} berada satu direktori dengan file \texttt{.tex}.}}%
            }%
        \end{center}
    \fi
}

\title{\textbf{Menggambar Grafik 2D dengan Euler Math Toolbox (EMT)}}
\subtitle{Panduan Komputasi dan Visualisasi Kurva Dua Dimensi pada Euler Math Toolbox}
\author{Catatan Praktikum dan Latihan Visualisasi 2D}
\date{\today}

\begin{document}

\frontmatter
\maketitle
\tableofcontents

\mainmatter
\begin{eulernotebook}

\eulerheading{Menggambar Grafik 2D dengan EMT}
\begin{eulercomment}
Notebook ini menjelaskan tentang cara menggambar berbagai kurva dan grafik 2D dengan perangkat lunak EMT. EMT menyediakan fungsi \texttt{plot2d()} untuk menggambar berbagai kurva dan grafik dua dimensi (2D).
\end{eulercomment}

\eulersubheading{Plot Dasar}
\begin{eulercomment}
Terdapat fungsi-fungsi plot yang sangat dasar. Terdapat koordinat layar, yang selalu memiliki rentang dari 0 hingga 1024 di setiap sumbunya, tidak peduli apakah layarnya persegi atau tidak. Selain itu, ada koordinat plot yang dapat diatur dengan \texttt{setplot()}. Pemetaan antara koordinat-koordinat tersebut bergantung pada jendela plot saat ini. Contohnya, \texttt{shrinkwindow()} bawaan menyisakan ruang untuk label sumbu dan judul plot.

Pada contoh ini, kita hanya menggambar beberapa garis acak dengan berbagai warna. Untuk detail mengenai fungsi-fungsi ini, pelajari fungsi-fungsi inti dari EMT.
\end{eulercomment}
\begin{eulerprompt}
>clg; // bersihkan layar
>window(0,0,1024,1024); // gunakan seluruh jendela
>setplot(0,1,0,1); // atur koordinat plot
>hold on; // mulai mode timpa (overwrite)
>n=100; X=random(n,2); Y=random(n,2); // dapatkan titik-titik acak
>colors=rgb(random(n),random(n),random(n)); // dapatkan warna acak
>loop 1 to n; color(colors[#]); plot(X[#],Y[#]); end; // gambar plot
>hold off; // akhiri mode timpa
>insimg; // sisipkan ke notebook
\end{eulerprompt}
\eulerimg{25}{images/Tugas 3-001.png}
\begin{eulerprompt}
>reset;
\end{eulerprompt}
\begin{eulercomment}
Grafik perlu ditahan (\texttt{hold}), karena perintah \texttt{plot()} akan membersihkan jendela plot.

Untuk menghapus semua yang telah kita buat, kita menggunakan \texttt{reset()}.

Untuk menampilkan gambar hasil plot di layar notebook, perintah \texttt{plot2d()} dapat diakhiri dengan titik dua (\texttt{:}). Cara lain adalah perintah \texttt{plot2d()} diakhiri dengan titik koma (\texttt{;}), kemudian menggunakan perintah \texttt{insimg()} untuk menampilkan gambar hasil plot.

Sebagai contoh lain, kita menggambar plot sebagai sisipan (\emph{inset}) di dalam plot lain. Hal ini dilakukan dengan mendefinisikan jendela plot yang lebih kecil. Perhatikan bahwa jendela ini tidak menyediakan ruang untuk label sumbu di luar jendela plot. Kita harus menambahkan sedikit margin untuk ini sesuai kebutuhan. Perhatikan bahwa kita menyimpan dan mengembalikan jendela penuh, dan menahan plot saat ini sementara kita menggambar sisipan tersebut.
\end{eulercomment}
\begin{eulerprompt}
>plot2d("x^3");
>xw=200; yw=100; ww=300; hw=300;
>ow=window();
>window(xw,yw,xw+ww,yw+hw);
>hold on;
>barclear(xw-50,yw-10,ww+60,ww+60);
>plot2d("x^4-x",grid=6):
\end{eulerprompt}
\eulerimg{25}{images/Tugas 3-002.png}
\begin{eulerprompt}
>hold off;
>window(ow);
\end{eulerprompt}
\begin{eulercomment}
Plot dengan banyak figur (gambar) dicapai dengan cara yang sama. Terdapat fungsi utilitas \texttt{figure()} untuk melakukan hal ini.
\end{eulercomment}

\eulersubheading{Aspek Plot}
\begin{eulercomment}
Plot bawaan menggunakan jendela plot persegi. Anda dapat mengubahnya dengan fungsi \texttt{aspect()}. Jangan lupa untuk mereset aspeknya nanti. Anda juga dapat mengubah nilai bawaan ini di menu dengan ``Set Aspect'' ke rasio aspek tertentu atau ke ukuran saat ini dari jendela grafis.

Namun Anda juga dapat mengubahnya hanya untuk satu plot. Untuk ini, ukuran saat ini dari area plot akan diubah, dan jendela diatur sedemikian rupa sehingga label memiliki ruang yang cukup.
\end{eulercomment}
\begin{eulerprompt}
>aspect(2); // rasio panjang dan lebar 2:1
>plot2d(["sin(x)","cos(x)"],0,2pi):
\end{eulerprompt}
\eulerimg{22}{images/Tugas 3-003.png}
\begin{eulerprompt}
>aspect();
>reset;
\end{eulerprompt}
\begin{eulercomment}
Fungsi \texttt{reset()} mengembalikan pengaturan bawaan plot termasuk rasio aspeknya.
\end{eulercomment}

\eulerheading{Plot 2D di Euler}
\begin{eulercomment}
Euler Math Toolbox (EMT) memiliki fitur plot dalam 2D, baik untuk data maupun fungsi. EMT menggunakan fungsi \texttt{plot2d}. Fungsi ini dapat menggambar fungsi dan data.

Dimungkinkan juga untuk melakukan plot di Maxima menggunakan Gnuplot atau di Python menggunakan Matplotlib.

Euler dapat menggambar plot 2D dari:
\begin{itemize}
    \item ekspresi,
    \item fungsi, variabel, atau kurva terparameterisasi,
    \item vektor nilai $x$-$y$,
    \item sekumpulan titik di bidang,
    \item kurva implisit dengan level atau rentang level,
    \item fungsi kompleks.
\end{itemize}

Gaya plot mencakup berbagai gaya untuk garis dan titik, plot batang, dan plot berbayang (\emph{shaded}).
\end{eulercomment}

\eulerheading{Plot Ekspresi atau Variabel}
\begin{eulercomment}
Sebuah ekspresi tunggal dalam \texttt{"x"} (misalnya \texttt{"4*x\textasciicircum{}2"}) atau nama sebuah fungsi (misalnya \texttt{"f"}) menghasilkan grafik dari fungsi tersebut.

Berikut adalah contoh paling dasar, yang menggunakan rentang bawaan dan mengatur rentang-$y$ yang tepat agar sesuai dengan plot fungsinya.

Catatan: Jika Anda mengakhiri baris perintah dengan titik dua \texttt{":"}, plot akan disisipkan ke dalam jendela teks. Jika tidak, tekan \texttt{TAB} untuk melihat plot jika jendela plot tertutup.
\end{eulercomment}
\begin{eulerprompt}
>plot2d("x^5"):
\end{eulerprompt}
\eulerimg{25}{images/Tugas 3-004.png}
\begin{eulerprompt}
>aspect(1.5); plot2d("x^3-x+x^2"):
\end{eulerprompt}
\eulerimg{30}{images/Tugas 3-005.png}
\begin{eulerprompt}
>a:=5.6; plot2d("exp(-a*x^2)/a"); insimg(30); // menampilkan gambar hasil plot setinggi 30 baris
\end{eulerprompt}
\eulerimg{29}{images/Tugas 3-006.png}
\begin{eulercomment}
Dari beberapa contoh sebelumnya Anda dapat melihat bahwa aslinya gambar plot menggunakan sumbu-$X$ dengan rentang nilai dari $-2$ sampai dengan $2$. Untuk mengubah rentang nilai $X$ dan $Y$, Anda dapat menambahkan nilai-nilai batas $X$ (dan $Y$) di belakang ekspresi yang digambar.

Rentang plot diatur dengan parameter-parameter yang ditetapkan berikut ini:
\begin{itemize}
    \item \texttt{a,b}: rentang-$x$ (bawaan $-2,2$)
    \item \texttt{c,d}: rentang-$y$ (bawaan: skala mengikuti nilai)
    \item \texttt{r}: sebagai alternatif, radius di sekitar pusat plot
    \item \texttt{cx,cy}: koordinat pusat plot (bawaan $0,0$)
\end{itemize}
\end{eulercomment}
\begin{eulerprompt}
>plot2d("x^3-x",-1,2):
\end{eulerprompt}
\eulerimg{30}{images/Tugas 3-007.png}
\begin{eulerprompt}
>plot2d("sin(x)",-2*pi,2*pi): // plot sin(x) pada interval [-2pi, 2pi]
\end{eulerprompt}
\eulerimg{30}{images/Tugas 3-008.png}
\begin{eulerprompt}
>plot2d("cos(x)","sin(3*x)",xmin=0,xmax=2pi):
\end{eulerprompt}
\eulerimg{30}{images/Tugas 3-009.png}
\begin{eulercomment}
Alternatif dari tanda titik dua adalah perintah \texttt{insimg(lines)}, yang menyisipkan plot dengan menempati jumlah baris teks yang ditentukan.

Pada opsi-opsi yang ada, plot dapat diatur agar muncul:
\begin{itemize}
    \item di jendela terpisah yang dapat diubah ukurannya,
    \item di jendela notebook.
\end{itemize}

Lebih banyak gaya dapat dicapai dengan perintah-perintah plot yang spesifik.

Bagaimanapun, tekan tombol tabulator untuk melihat plot, jika plot tersebut tersembunyi.

Untuk membagi jendela menjadi beberapa plot, gunakan perintah \texttt{figure()}. Pada contoh ini, kita menggambar $x^1$ hingga $x^4$ ke dalam 4 bagian jendela. \texttt{figure(0)} mereset ke jendela bawaan.
\end{eulercomment}
\begin{eulerprompt}
>reset;
>figure(2,2); ...
>for n=1 to 4; figure(n); plot2d("x^"+n); end; ...
>figure(0):
\end{eulerprompt}
\eulerimg{25}{images/Tugas 3-010.png}
\begin{eulercomment}
Di \texttt{plot2d()}, terdapat gaya alternatif yang tersedia dengan \texttt{grid=x}. Sebagai gambaran, kita menampilkan berbagai gaya grid dalam satu figur (lihat di bawah untuk perintah \texttt{figure()}). Gaya \texttt{grid=0} tidak disertakan karena gaya tersebut tidak menampilkan grid dan bingkai.
\end{eulercomment}
\begin{eulerprompt}
>figure(3,3); ...
>for k=1:9; figure(k); plot2d("x^3-x",-2,1,grid=k); end; ...
>figure(0):
\end{eulerprompt}
\eulerimg{25}{images/Tugas 3-011.png}
\begin{eulercomment}
Jika argumen pada \texttt{plot2d()} adalah ekspresi yang diikuti oleh empat angka, angka-angka ini merupakan rentang $x$ dan $y$ untuk plot tersebut.

Secara alternatif, \texttt{a, b, c, d} dapat ditentukan sebagai parameter yang ditetapkan seperti \texttt{a=...} dst.

Pada contoh berikut, kita mengubah gaya grid, menambahkan label, dan menggunakan label vertikal untuk sumbu-$y$.
\end{eulercomment}
\begin{eulerprompt}
>aspect(1.5); plot2d("sin(x)",0,2pi,-1.2,1.2,grid=3,xl="x",yl="sin(x)"):
\end{eulerprompt}
\eulerimg{30}{images/Tugas 3-012.png}
\begin{eulerprompt}
>plot2d("sin(x)+cos(2*x)",0,4pi):
\end{eulerprompt}
\eulerimg{30}{images/Tugas 3-013.png}
\begin{eulercomment}
Gambar yang dihasilkan dengan menyisipkan plot ke dalam jendela teks disimpan di direktori yang sama dengan notebook, secara bawaan berada di subdirektori bernama \texttt{"images"}. Gambar-gambar tersebut juga digunakan oleh ekspor HTML.

Anda cukup menandai gambar apa pun dan menyalinnya ke papan klip (\emph{clipboard}) dengan \texttt{Ctrl-C}. Tentu saja, Anda juga dapat mengekspor grafik saat ini dengan fungsi-fungsi di menu File.

Fungsi atau ekspresi di \texttt{plot2d} dievaluasi secara adaptif. Untuk kecepatan lebih, matikan plot adaptif dengan \texttt{<adaptive} dan tentukan jumlah subinterval dengan \texttt{n=...}. Ini hanya diperlukan pada kasus-kasus yang jarang terjadi.
\end{eulercomment}
\begin{eulerprompt}
>plot2d("sign(x)*exp(-x^2)",-1,1,<adaptive,n=10000):
\end{eulerprompt}
\eulerimg{30}{images/Tugas 3-014.png}
\begin{eulerprompt}
>plot2d("x^x",r=1.2,cx=1,cy=1):
\end{eulerprompt}
\eulerimg{30}{images/Tugas 3-015.png}
\begin{eulercomment}
Perhatikan bahwa $x^x$ tidak terdefinisi untuk $x \le 0$. Fungsi \texttt{plot2d} menangkap galat ini, dan mulai menggambar segera setelah fungsinya terdefinisi. Ini berlaku untuk semua fungsi yang mengembalikan \texttt{NAN} di luar rentang definisinya.
\end{eulercomment}
\begin{eulerprompt}
>plot2d("log(x)",-0.1,2):
\end{eulerprompt}
\eulerimg{30}{images/Tugas 3-016.png}
\begin{eulercomment}
Parameter \texttt{square=true} (atau \texttt{>square}) memilih rentang-$y$ secara otomatis sehingga hasilnya adalah jendela plot berbentuk persegi. Perhatikan bahwa secara bawaan, Euler menggunakan ruang persegi di dalam jendela plot.
\end{eulercomment}
\begin{eulerprompt}
>plot2d("x^3-x",>square):
\end{eulerprompt}
\eulerimg{30}{images/Tugas 3-017.png}
\begin{eulerprompt}
>plot2d(''integrate("sin(x)*exp(-x^2)",0,x)'',0,2): // gambar plot integral
\end{eulerprompt}
\eulerimg{30}{images/Tugas 3-018.png}
\begin{eulercomment}
Jika Anda membutuhkan lebih banyak ruang untuk label-$y$, panggil \texttt{shrinkwindow()} dengan parameter yang lebih kecil, atau tetapkan nilai positif untuk \texttt{"smaller"} di \texttt{plot2d()}.
\end{eulercomment}
\begin{eulerprompt}
>plot2d("gamma(x)",1,10,yl="y-values",smaller=6,<vertical):
\end{eulerprompt}
\eulerimg{30}{images/Tugas 3-019.png}
\begin{eulercomment}
Ekspresi simbolik juga dapat digunakan, karena disimpan sebagai ekspresi string sederhana.
\end{eulercomment}
\begin{eulerprompt}
>x=linspace(0,2pi,1000); plot2d(sin(5x),cos(7x)):
\end{eulerprompt}
\eulerimg{30}{images/Tugas 3-020.png}
\begin{eulerprompt}
>a:=5.6; expr &= exp(-a*x^2)/a; // definisikan ekspresi
>plot2d(expr,-2,2): // gambar plot dari -2 hingga 2
\end{eulerprompt}
\eulerimg{30}{images/Tugas 3-021.png}
\begin{eulerprompt}
>plot2d(expr,r=1,thickness=2): // gambar plot dalam persegi di sekitar (0,0)
\end{eulerprompt}
\eulerimg{30}{images/Tugas 3-022.png}
\begin{eulerprompt}
>plot2d(&diff(expr,x),>add,style="--",color=red): // tambahkan plot lain
\end{eulerprompt}
\eulerimg{30}{images/Tugas 3-023.png}
\begin{eulerprompt}
>plot2d(&diff(expr,x,2),a=-2,b=2,c=-2,d=1): // gambar plot dalam persegi panjang
\end{eulerprompt}
\eulerimg{30}{images/Tugas 3-024.png}
\begin{eulerprompt}
>plot2d(&diff(expr,x),a=-2,b=2,>square): // pertahankan plot berbentuk persegi
\end{eulerprompt}
\eulerimg{30}{images/Tugas 3-025.png}
\begin{eulerprompt}
>plot2d("x^2",0,1,steps=1,color=red,n=10):
\end{eulerprompt}
\eulerimg{30}{images/Tugas 3-026.png}
\begin{eulerprompt}
>plot2d("x^2",>add,steps=2,color=blue,n=10):
\end{eulerprompt}
\eulerimg{30}{images/Tugas 3-027.png}

\eulerheading{Fungsi dengan satu Parameter}
\begin{eulercomment}
Fungsi plot terpenting untuk plot planar adalah \texttt{plot2d()}. Fungsi ini diimplementasikan dalam bahasa Euler di file \texttt{"plot.e"}, yang dimuat pada saat memulai program.

Berikut adalah beberapa contoh penggunaan fungsi. Seperti biasa di EMT, untuk fungsi-fungsi yang bekerja pada fungsi atau ekspresi lain, Anda dapat memberikan parameter tambahan (selain $x$) yang bukan merupakan variabel global ke fungsi tersebut dengan parameter titik koma atau dengan kumpulan pemanggilan (\emph{call collection}).
\end{eulercomment}
\begin{eulerprompt}
>plot2d("x^2-4*x+3", -5, 5):
\end{eulerprompt}
\eulerimg{30}{images/Tugas 3-028.png}
\begin{eulerprompt}
>plot2d("x^5-10*x^3+6*x^2-1100", -75, 75):
\end{eulerprompt}
\eulerimg{30}{images/Tugas 3-029.png}
\begin{eulerprompt}
>function h(x)&= cos(x^2-1);
>plot2d("h", -10, 10):
\end{eulerprompt}
\eulerimg{30}{images/Tugas 3-030.png}
\begin{eulerprompt}
>plot2d("h", -5, 5):
\end{eulerprompt}
\eulerimg{30}{images/Tugas 3-031.png}
\begin{eulerprompt}
>plot2d("h", -100, 100):
\end{eulerprompt}
\eulerimg{30}{images/Tugas 3-032.png}
\begin{eulerprompt}
>function r(y)&=cos(6*y);
>plot2d("r",-10,10):
\end{eulerprompt}
\eulerimg{30}{images/Tugas 3-033.png}
\begin{eulerprompt}
>function f(x,a) := x^2/a+a*x^2-x; // definisikan sebuah fungsi
\end{eulerprompt}
\begin{eulercomment}
Berikut adalah ringkasan dari fungsi-fungsi yang diterima:
\begin{itemize}
    \item ekspresi atau ekspresi simbolik dalam $x$,
    \item fungsi atau fungsi simbolik berdasarkan nama seperti \texttt{"f"},
    \item fungsi simbolik hanya dengan nama \texttt{f}.
\end{itemize}

Fungsi \texttt{plot2d()} juga menerima fungsi simbolik. Untuk fungsi simbolik, namanya saja sudah berfungsi.
\end{eulercomment}
\begin{eulerprompt}
>function f(x) &= diff(x^x,x)
\end{eulerprompt}
\begin{euleroutput}
        x
       x  (log(x) + 1)
\end{euleroutput}
\begin{eulerprompt}
>plot2d(f,0,2):
\end{eulerprompt}
\eulerimg{30}{images/Tugas 3-034.png}
\begin{eulercomment}
Tentu saja, untuk ekspresi atau ekspresi simbolik, nama variabelnya sudah cukup untuk menggambarnya.
\end{eulercomment}
\begin{eulerprompt}
>expr &= sin(x)*exp(-x)
\end{eulerprompt}
\begin{euleroutput}
        - x
       E    sin(x)
\end{euleroutput}
\begin{eulerprompt}
>plot2d(expr,0,3pi):
\end{eulerprompt}
\eulerimg{30}{images/Tugas 3-035.png}
\begin{eulerprompt}
>function f(x) &= x^x;
>plot2d(f,r=1,cx=1,cy=1,color=blue,thickness=2);
>plot2d(&diff(f(x),x),>add,color=red,style="-.-"):
\end{eulerprompt}
\eulerimg{30}{images/Tugas 3-036.png}
\begin{eulercomment}
Untuk gaya garis terdapat berbagai opsi:
\begin{itemize}
    \item \texttt{style="..."}: Pilih dari \texttt{"-"}, \texttt{"--"}, \texttt{"-."}, \texttt{"."}, \texttt{".-."}, \texttt{"-.-"}.
    \item \texttt{color}: Lihat di bawah untuk warna.
    \item \texttt{thickness}: Nilai bawaan adalah 1.
\end{itemize}

Warna dapat dipilih sebagai salah satu warna bawaan, atau sebagai warna RGB:
\begin{itemize}
    \item \texttt{0..15}: indeks warna bawaan.
    \item konstanta warna: \texttt{white}, \texttt{black}, \texttt{red}, \texttt{green}, \texttt{blue}, \texttt{cyan}, \texttt{olive}, \texttt{lightgray}, \texttt{gray}, \texttt{darkgray}, \texttt{orange}, \texttt{lightgreen}, \texttt{turquoise}, \texttt{lightblue}, \texttt{lightorange}, \texttt{yellow}.
    \item \texttt{rgb(red,green,blue)}: parameternya adalah bilangan real di $[0,1]$.
\end{itemize}
\end{eulercomment}
\begin{eulerprompt}
>plot2d("exp(-x^2)",r=2,color=red,thickness=3,style="--"):
\end{eulerprompt}
\eulerimg{30}{images/Tugas 3-037.png}
\begin{eulercomment}
Berikut adalah tampilan warna-warna yang telah ditentukan sebelumnya di EMT.
\end{eulercomment}
\begin{eulerprompt}
>aspect(2); columnsplot(ones(1,16),lab=0:15,grid=0,color=0:15):
\end{eulerprompt}
\eulerimg{22}{images/Tugas 3-038.png}
\begin{eulercomment}
Namun Anda dapat menggunakan warna apa saja.
\end{eulercomment}
\begin{eulerprompt}
>columnsplot(ones(1,16),grid=0,color=rgb(0,0,linspace(0,1,15))):
\end{eulerprompt}
\eulerimg{22}{images/Tugas 3-039.png}

\eulerheading{Menggambar Beberapa Kurva pada bidang koordinat yang sama}
\begin{eulercomment}
Menggambar lebih dari satu fungsi (fungsi ganda) ke dalam satu jendela dapat dilakukan dengan berbagai cara. Salah satu metodenya adalah menggunakan \texttt{>add} untuk beberapa pemanggilan \texttt{plot2d} di seluruh pemanggilan, kecuali pemanggilan pertama. Kita sudah menggunakan fitur ini pada contoh-contoh di atas.
\end{eulercomment}
\begin{eulerprompt}
>aspect(); plot2d("cos(x)",r=2,grid=6); plot2d("x",style=".",>add):
\end{eulerprompt}
\eulerimg{25}{images/Tugas 3-040.png}
\begin{eulerprompt}
>aspect(1.5); plot2d("sin(x)",0,2pi); plot2d("cos(x)",color=blue,style="--",>add):
\end{eulerprompt}
\eulerimg{30}{images/Tugas 3-041.png}
\begin{eulercomment}
Salah satu kegunaan \texttt{>add} adalah untuk menambahkan titik pada kurva.
\end{eulercomment}
\begin{eulerprompt}
>plot2d("sin(x)",0,pi); plot2d(2,sin(2),>points,>add):
\end{eulerprompt}
\eulerimg{30}{images/Tugas 3-042.png}
\begin{eulercomment}
Kita menambahkan titik potong beserta labelnya (pada posisi \texttt{"cl"} untuk kiri tengah/\emph{center left}), dan menyisipkan hasilnya ke dalam notebook. Kita juga menambahkan judul pada plot tersebut.
\end{eulercomment}
\begin{eulerprompt}
>plot2d(["cos(x)","x"],r=1.1,cx=0.5,cy=0.5, ...
>  color=[black,blue],style=["-","."], ...
>  grid=1);
>x0=solve("cos(x)-x",1); ...
>  plot2d(x0,x0,>points,>add,title="Intersection Demo"); ...
>  label("cos(x) = x",x0,x0,pos="cl",offset=20):
\end{eulerprompt}
\eulerimg{30}{images/Tugas 3-043.png}
\begin{eulercomment}
Pada demo berikut, kita menggambar fungsi $\mathrm{sinc}(x)=\sin(x)/x$ beserta ekspansi Taylor ke-8 dan ke-16 nya. Kita menghitung ekspansi ini menggunakan Maxima melalui ekspresi simbolik.
Plot ini dilakukan pada perintah multi-baris berikut dengan tiga pemanggilan \texttt{plot2d()}. Pemanggilan kedua dan ketiga memiliki bendera (\emph{flag}) \texttt{>add} yang disetel, yang membuat plot-plot tersebut menggunakan rentang dari plot sebelumnya.

Kita menambahkan kotak label untuk menjelaskan fungsi-fungsi tersebut.
\end{eulercomment}
\begin{eulerprompt}
>$taylor(sin(x)/x,x,0,4)
\end{eulerprompt}
\begin{eulerformula}
\[
\frac{x^4}{120}-\frac{x^2}{6}+1
\]
\end{eulerformula}
\begin{eulerprompt}
>plot2d("sinc(x)",0,4pi,color=green,thickness=2); ...
>  plot2d(&taylor(sin(x)/x,x,0,8),>add,color=blue,style="--"); ...
>  plot2d(&taylor(sin(x)/x,x,0,16),>add,color=red,style="-.-"); ...
>  labelbox(["sinc","T8","T16"],styles=["-","--","-.-"], ...
>    colors=[black,blue,red]):
\end{eulerprompt}
\eulerimg{30}{images/Tugas 3-045.png}
\begin{eulercomment}
Pada contoh berikut, kita menghasilkan Polinomial Bernstein.
\end{eulercomment}
\begin{eulerformula}
\[
B_i(x) = \binom{n}{i} x^i (1-x)^{n-i}
\]
\end{eulerformula}
\begin{eulerprompt}
>plot2d("(1-x)^10",0,1); // gambar fungsi pertama
>for i=1 to 10; plot2d("bin(10,i)*x^i*(1-x)^(10-i)",>add); end;
>insimg;
\end{eulerprompt}
\eulerimg{30}{images/Tugas 3-047.png}
\begin{eulercomment}
Metode kedua adalah menggunakan sepasang matriks nilai-$x$ dan matriks nilai-$y$ dengan ukuran yang sama.

Kita menghasilkan matriks nilai dengan satu Polinomial Bernstein di setiap barisnya. Untuk ini, kita cukup menggunakan vektor kolom \texttt{k}. Silakan lihat bagian pengenalan tentang bahasa matriks untuk mempelajari detail lebih lanjut.
\end{eulercomment}
\begin{eulerprompt}
>x=linspace(0,1,500);
>n=10; k=(0:n)'; // n adalah skalar/vektor baris, k adalah vektor kolom
>y=bin(n,k)*x^k*(1-x)^(n-k); // maka y adalah sebuah matriks
>plot2d(x,y):
\end{eulerprompt}
\eulerimg{30}{images/Tugas 3-048.png}
\begin{eulercomment}
Perhatikan bahwa parameter warna dapat berupa vektor. Maka setiap warna digunakan untuk setiap baris dari matriks tersebut.
\end{eulercomment}
\begin{eulerprompt}
>x=linspace(0,1,200); y=x^(1:10)'; plot2d(x,y,color=1:10):
\end{eulerprompt}
\eulerimg{30}{images/Tugas 3-049.png}
\begin{eulercomment}
Metode lainnya adalah menggunakan vektor ekspresi (string). Anda kemudian dapat menggunakan larik (\emph{array}) warna, larik gaya, dan larik ketebalan dengan panjang yang sama.
\end{eulercomment}
\begin{eulerprompt}
>plot2d(["sin(x)","cos(x)"],0,2pi,color=4:5): 
\end{eulerprompt}
\eulerimg{30}{images/Tugas 3-050.png}
\begin{eulerprompt}
>plot2d(["sin(x)","cos(x)"],0,2pi): // gambar vektor ekspresi
\end{eulerprompt}
\eulerimg{30}{images/Tugas 3-051.png}
\begin{eulercomment}
Kita dapat memperoleh vektor semacam itu dari Maxima menggunakan \texttt{makelist()} dan \texttt{mxm2str()}.
\end{eulercomment}
\begin{eulerprompt}
>v &= makelist(binomial(10,i)*x^i*(1-x)^(10-i),i,0,10) // buat list
\end{eulerprompt}
\begin{euleroutput}
         10            9              8  2             7  3             6  4
[(1 - x)  , 10 (1 - x)  x, 45 (1 - x)  x , 120 (1 - x)  x , 210 (1 - x)  x , 
          5  5             4  6             3  7            2  8              9   10
252 (1 - x)  x , 210 (1 - x)  x , 120 (1 - x)  x , 45 (1 - x)  x , 10 (1 - x) x , x  ]
\end{euleroutput}
\begin{eulerprompt}
>mxm2str(v) // dapatkan vektor string dari vektor simbolik
\end{eulerprompt}
\begin{euleroutput}
(1-x)^10
10*(1-x)^9*x
45*(1-x)^8*x^2
120*(1-x)^7*x^3
210*(1-x)^6*x^4
252*(1-x)^5*x^5
210*(1-x)^4*x^6
120*(1-x)^3*x^7
45*(1-x)^2*x^8
10*(1-x)*x^9
x^10
\end{euleroutput}
\begin{eulerprompt}
>plot2d(mxm2str(v),0,1): // gambar fungsi-fungsi
\end{eulerprompt}
\eulerimg{30}{images/Tugas 3-052.png}
\begin{eulercomment}
Alternatif lain adalah menggunakan bahasa matriks Euler.

Jika sebuah ekspresi menghasilkan matriks fungsi, dengan satu fungsi di setiap baris, semua fungsi ini akan digambar ke dalam satu plot.

Untuk ini, gunakan vektor parameter dalam bentuk vektor kolom. Jika larik warna ditambahkan, warna tersebut akan digunakan untuk setiap baris plot.
\end{eulercomment}
\begin{eulerprompt}
>n=(1:10)'; plot2d("x^n",0,1,color=1:10):
\end{eulerprompt}
\eulerimg{30}{images/Tugas 3-053.png}
\begin{eulercomment}
Ekspresi dan fungsi satu baris dapat melihat variabel global.

Jika Anda tidak dapat menggunakan variabel global, Anda harus menggunakan fungsi dengan parameter ekstra, dan memberikan parameter ini sebagai parameter titik koma.

Hati-hati, pastikan menaruh semua parameter yang ditetapkan ke bagian akhir perintah \texttt{plot2d}. Pada contoh ini kita memberikan \texttt{a=5} ke fungsi \texttt{f}, yang kita gambar dari $-10$ hingga $10$.
\end{eulercomment}
\begin{eulerprompt}
>function f(x,a) := 1/a*exp(-x^2/a); ...
>plot2d("f",-10,10;5,thickness=2,title="a=5"):
\end{eulerprompt}
\eulerimg{30}{images/Tugas 3-054.png}
\begin{eulercomment}
Sebagai alternatif, gunakan koleksi dengan nama fungsi dan semua parameter ekstra. Daftar khusus ini disebut kumpulan pemanggilan (\emph{call collections}), dan itu adalah cara yang disukai untuk meneruskan argumen ke sebuah fungsi yang mana fungsi itu sendiri diteruskan sebagai argumen ke fungsi lain.

Pada contoh berikut, kita menggunakan perulangan (\emph{loop}) untuk menggambar beberapa fungsi (lihat tutorial tentang pemrograman perulangan \texttt{for}).
\end{eulercomment}
\begin{eulerprompt}
>plot2d({{"f",1}},-10,10); ...
>for a=2:10; plot2d({{"f",a}},>add); end:
\end{eulerprompt}
\eulerimg{30}{images/Tugas 3-055.png}
\begin{eulercomment}
Kita bisa mencapai hasil yang sama dengan cara berikut menggunakan bahasa matriks EMT. Setiap baris dari matriks \texttt{f(x,a)} adalah satu fungsi. Selain itu, kita dapat mengatur warna untuk setiap baris matriks tersebut. Klik dua kali pada fungsi \texttt{getspectral()} untuk melihat penjelasannya.
\end{eulercomment}
\begin{eulerprompt}
>x=-10:0.01:10; a=(1:10)'; plot2d(x,f(x,a),color=getspectral(a/10)):
\end{eulerprompt}
\eulerimg{30}{images/Tugas 3-056.png}

\eulersubheading{Label Teks}
\begin{eulercomment}
Dekorasi sederhana dapat berupa:
\begin{itemize}
    \item judul dengan \texttt{title="..."}
    \item label $x$ dan $y$ dengan \texttt{xl="..."}, \texttt{yl="..."}
    \item label teks lainnya dengan \texttt{label("...",x,y)}
\end{itemize}

Perintah \texttt{label} akan menggambar di plot saat ini pada koordinat plot $(x,y)$. Perintah ini dapat mengambil argumen posisional.
\end{eulercomment}
\begin{eulerprompt}
>plot2d("x^3-x",-1,2,title="y=x^3-x",yl="y",xl="x"):
\end{eulerprompt}
\eulerimg{30}{images/Tugas 3-057.png}
\begin{eulerprompt}
>expr := "log(x)/x"; ...
>  plot2d(expr,0.5,5,title="y="+expr,xl="x",yl="y"); ...
>  label("(1,0)",1,0); label("Max",E,expr(E),pos="lc"):
\end{eulerprompt}
\eulerimg{30}{images/Tugas 3-058.png}
\begin{eulercomment}
Terdapat juga fungsi \texttt{labelbox()}, yang dapat menampilkan fungsi-fungsi dan teks. Fungsi ini mengambil vektor string dan warna, satu item untuk setiap fungsi.
\end{eulercomment}
\begin{eulerprompt}
>function f(x) &= x^2*exp(-x^2); ...
>plot2d(&f(x),a=-3,b=3,c=-1,d=1); ...
>plot2d(&diff(f(x),x),>add,color=blue,style="--"); ...
>labelbox(["function","derivative"],styles=["-","--"], ...
>   colors=[black,blue],w=0.4):
\end{eulerprompt}
\eulerimg{30}{images/Tugas 3-059.png}
\begin{eulercomment}
Kotak ini ditambatkan di kanan atas secara bawaan, namun \texttt{>left} menambatkannya di kiri atas. Anda dapat memindahkannya ke tempat mana pun yang Anda suka. Posisi tambatan adalah sudut kanan atas dari kotak, dan angka-angkanya adalah pecahan dari ukuran jendela grafis. Lebarnya otomatis.

Untuk plot titik, kotak label juga berfungsi. Tambahkan parameter \texttt{>points}, atau vektor bendera (\emph{flags}), satu untuk setiap label.

Pada contoh berikut, hanya terdapat satu fungsi. Jadi kita dapat menggunakan string daripada vektor string. Kita mengatur warna teks menjadi hitam untuk contoh ini.
\end{eulercomment}
\begin{eulerprompt}
>n=10; plot2d(0:n,bin(n,0:n),>addpoints); ...
>labelbox("Binomials",styles="[]",>points,x=0.1,y=0.1, ...
>tcolor=black,>left):
\end{eulerprompt}
\eulerimg{30}{images/Tugas 3-060.png}
\begin{eulercomment}
Gaya plot ini juga tersedia di \texttt{statplot()}. Seperti di \texttt{plot2d()}, warna dapat diatur untuk setiap baris plot. Ada lebih banyak plot khusus untuk tujuan statistik (lihat tutorial tentang statistik).
\end{eulercomment}
\begin{eulerprompt}
>statplot(1:10,random(2,10),color=[red,blue]):
\end{eulerprompt}
\eulerimg{30}{images/Tugas 3-061.png}
\begin{eulercomment}
Fitur yang serupa adalah fungsi \texttt{textbox()}.

Lebarnya secara bawaan adalah lebar maksimal dari baris-baris teks. Tetapi ia juga dapat diatur oleh pengguna.
\end{eulercomment}
\begin{eulerprompt}
>function f(x) &= exp(-x)*sin(2*pi*x); ...
>plot2d("f(x)",0,2pi); ...
>textbox(latex("\text{Example of a damped oscillation}\ f(x)=e^{-x}sin(2\pi x)"),w=0.85):
\end{eulerprompt}
\eulerimg{30}{images/Tugas 3-062.png}
\begin{eulercomment}
Label teks, judul, kotak label, dan teks lainnya dapat berisi string Unicode (lihat sintaks EMT untuk info lebih lanjut tentang string Unicode).
\end{eulercomment}
\begin{eulerprompt}
>plot2d("x^3-x",title=u"x &rarr; x&sup3; - x"):
\end{eulerprompt}
\eulerimg{30}{images/Tugas 3-063.png}
\begin{eulercomment}
Label pada sumbu-$x$ dan sumbu-$y$ dapat dibuat vertikal, begitu pula sumbunya.
\end{eulercomment}
\begin{eulerprompt}
>plot2d("sinc(x)",0,2pi,xl="x",yl=u"x &rarr; sinc(x)",>vertical):
\end{eulerprompt}
\eulerimg{30}{images/Tugas 3-064.png}

\eulersubheading{LaTeX}
\begin{eulercomment}
Anda juga dapat menggambar formula LaTeX jika Anda telah menginstal sistem LaTeX (misalnya MiKTeX). Jalur (\emph{path}) ke file biner \texttt{"latex"} dan \texttt{"dvipng"} harus berada di dalam \emph{path} sistem, atau Anda harus mengatur LaTeX di menu opsi.

Perhatikan bahwa mem-parsing LaTeX cukup lambat. Jika Anda ingin menggunakan LaTeX pada plot animasi, Anda sebaiknya memanggil \texttt{latex()} sebelum perulangan sekali saja dan menggunakan hasilnya (sebuah gambar di dalam matriks RGB).

Pada plot berikut, kita menggunakan LaTeX untuk label $x$ dan $y$, sebuah label, kotak label, dan judul plotnya.
\end{eulercomment}
\begin{eulerprompt}
>plot2d("exp(-x)*sin(x)/x",a=0,b=2pi,c=0,d=1,grid=6,color=blue, ...
>  title=latex("\text{Function $\Phi$}"), ...
>  xl=latex("\phi"),yl=latex("\Phi(\phi)")); ...
>textbox( ...
>  latex("\Phi(\phi) = e^{-\phi} \frac{\sin(\phi)}{\phi}"),x=0.8,y=0.5); ...
>label(latex("\Phi",color=blue),1,0.4):
\end{eulerprompt}
\eulerimg{30}{images/Tugas 3-065.png}
\begin{eulercomment}
Sering kali kita menginginkan spasi dan label teks yang tidak seragam pada sumbu-$x$. Kita dapat menggunakan \texttt{xaxis()} dan \texttt{yaxis()} seperti yang akan kita tunjukkan nanti.

Cara termudah adalah membuat plot kosong dengan bingkai menggunakan \texttt{grid=4}, dan kemudian menambahkan grid dengan \texttt{ygrid()} dan \texttt{xgrid()}. Pada contoh berikut, kita menggunakan tiga string LaTeX untuk label pada sumbu-$x$ dengan \texttt{xtick()}.
\end{eulercomment}
\begin{eulerprompt}
>plot2d("sinc(x)",0,2pi,grid=4,<ticks); ...
>ygrid(-2:0.5:2,grid=6); ...
>xgrid([0:2]*pi,<ticks,grid=6); ...
>xtick([0,pi,2pi],["0","\pi","2\pi"],>latex):
\end{eulerprompt}
\eulerimg{30}{images/Tugas 3-066.png}
\begin{eulercomment}
Tentu saja, fungsi juga dapat digunakan.
\end{eulercomment}
\begin{eulerprompt}
>function map f(x) ...
\end{eulerprompt}
\begin{eulerudf}
if x>0 then return x^4
else return x^2
endif
endfunction
\end{eulerudf}
\begin{eulercomment}
Parameter \texttt{"map"} membantu dalam menggunakan fungsi tersebut untuk vektor. Untuk plotnya, ini tidak diperlukan. Namun untuk mendemonstrasikan bahwa vektorisasi itu berguna, kita menambahkan beberapa titik kunci ke plot pada $x=-1$, $x=0$, dan $x=1$.

Pada plot berikut, kita juga memasukkan beberapa kode LaTeX. Kita menggunakannya untuk dua label dan sebuah kotak teks. Tentu saja, Anda hanya dapat menggunakan LaTeX jika Anda telah menginstal LaTeX dengan benar.
\end{eulercomment}
\begin{eulerprompt}
>plot2d("f",-1,1,xl="x",yl="f(x)",grid=6); ...
>plot2d([-1,0,1],f([-1,0,1]),>points,>add); ...
>label(latex("x^3"),0.72,f(0.72)); ...
>label(latex("x^2"),-0.52,f(-0.52),pos="ll"); ...
>textbox( ...
>  latex("f(x)=\begin{cases} x^3 & x>0 \\ x^2 & x \le 0\end{cases}"), ...
>  x=0.7,y=0.2):
\end{eulerprompt}
\eulerimg{30}{images/Tugas 3-067.png}

\eulersubheading{Interaksi Pengguna}
\begin{eulercomment}
Saat menggambar sebuah fungsi atau ekspresi, parameter \texttt{>user} memungkinkan pengguna untuk memperbesar (\emph{zoom}) dan menggeser plot menggunakan tombol kursor atau mouse. Pengguna dapat:
\begin{itemize}
    \item memperbesar dengan \texttt{+} atau \texttt{-},
    \item menggeser plot dengan tombol kursor,
    \item memilih jendela plot dengan mouse,
    \item mereset tampilan dengan tombol spasi,
    \item keluar dengan tombol \texttt{Return/Enter}.
\end{itemize}

Tombol spasi akan mereset plot ke jendela plot semula.

Saat menggambar data, bendera \texttt{>user} hanya akan menunggu ketukan tombol.
\end{eulercomment}
\begin{eulerprompt}
>plot2d({{"x^3-a*x",a=1}},>user,title="Press any key!"):
\end{eulerprompt}
\eulerimg{30}{images/Tugas 3-068.png}
\begin{eulerprompt}
>plot2d("exp(x)*sin(x)",user=true, ...
>  title="+/- or cursor keys (return to exit)"):
\end{eulerprompt}
\eulerimg{30}{images/Tugas 3-069.png}
\begin{eulercomment}
Hal berikut mendemonstrasikan cara lanjutan dari interaksi pengguna (lihat tutorial tentang pemrograman untuk detailnya).

Fungsi bawaan \texttt{mousedrag()} menunggu peristiwa mouse atau keyboard. Fungsi ini melaporkan mouse ditekan (\emph{down}), mouse digerakkan (\emph{moved}), atau mouse dilepas (\emph{up}), serta penekanan tombol. Fungsi \texttt{dragpoints()} memanfaatkan ini, dan memungkinkan pengguna menyeret (\emph{drag}) titik mana pun di dalam plot.

Pertama-tama kita memerlukan sebuah fungsi plot. Sebagai contoh, kita melakukan interpolasi pada 5 titik dengan polinomial. Fungsi ini harus menggambar di area plot yang tetap.
\end{eulercomment}
\begin{eulerprompt}
>function plotf(xp,yp,select) ...
\end{eulerprompt}
\begin{eulerudf}
d=interp(xp,yp);
plot2d("interpval(xp,d,x)";d,xp,r=2);
plot2d(xp,yp,>points,>add);
if select>0 then
    plot2d(xp[select],yp[select],color=red,>points,>add);
endif;
title("Drag one point, or press space or return!");
endfunction
\end{eulerudf}
\begin{eulercomment}
Perhatikan parameter titik koma di \texttt{plot2d} (\texttt{d} dan \texttt{xp}), yang diteruskan ke evaluasi fungsi \texttt{interp()}. Tanpa ini, kita harus menulis fungsi \texttt{plotinterp()} terlebih dahulu, yang mengakses nilai-nilai secara global.

Sekarang kita membuat beberapa nilai acak, dan membiarkan pengguna menyeret titik-titiknya.
\end{eulercomment}
\begin{eulerprompt}
>t=-1:0.5:1; dragpoints("plotf",t,random(size(t))-0.5):
\end{eulerprompt}
\eulerimg{30}{images/Tugas 3-070.png}
\begin{eulercomment}
Ada juga fungsi yang menggambar fungsi lain bergantung pada vektor parameter, dan memungkinkan pengguna menyesuaikan parameter-parameter ini.

Pertama kita membutuhkan fungsi plotnya.
\end{eulercomment}
\begin{eulerprompt}
>function plotf([a,b]) := plot2d("exp(a*x)*cos(2pi*b*x)",0,2pi;a,b);
\end{eulerprompt}
\begin{eulercomment}
Kemudian kita membutuhkan nama untuk parameternya, nilai awal, dan matriks $n \times 2$ untuk rentangnya, serta baris judul opsional.
Ada penggeser (\emph{slider}) interaktif yang nilainya dapat diatur oleh pengguna. Fungsi \texttt{dragvalues()} menyediakan hal ini.
\end{eulercomment}
\begin{eulerprompt}
>dragvalues("plotf",["a","b"],[-1,2],[[-2,2];[1,10]], ...
>  heading="Drag these values:",hcolor=black):
\end{eulerprompt}
\eulerimg{30}{images/Tugas 3-071.png}
\begin{eulercomment}
Nilai yang diseret dimungkinkan untuk dibatasi pada bilangan bulat. Sebagai contoh, kita menulis sebuah fungsi plot yang menggambar polinomial Taylor derajat $n$ pada fungsi kosinus.
\end{eulercomment}
\begin{eulerprompt}
>function plotf(n) ...
\end{eulerprompt}
\begin{eulerudf}
plot2d("cos(x)",0,2pi,>square,grid=6);
plot2d(&"taylor(cos(x),x,0,@n)",color=blue,>add);
textbox("Taylor polynomial of degree "+n,0.1,0.02,style="t",>left);
endfunction
\end{eulerudf}
\begin{eulercomment}
Sekarang kita membiarkan derajat $n$ bervariasi dari 0 hingga 20 dalam 20 pemberhentian. Hasil dari \texttt{dragvalues()} digunakan untuk menggambar sketsa dengan $n$ ini, dan untuk menyisipkan plot ke dalam notebook.
\end{eulercomment}
\begin{eulerprompt}
>nd=dragvalues("plotf","degree",2,[0,20],20,y=0.8, ...
>   heading="Drag the value:"); ...
>plotf(nd):
\end{eulerprompt}
\eulerimg{30}{images/Tugas 3-072.png}
\begin{eulercomment}
Berikut adalah demonstrasi sederhana dari fungsi tersebut. Pengguna dapat menggambar di atas jendela plot, meninggalkan jejak titik-titik.
\end{eulercomment}
\begin{eulerprompt}
>function dragtest ...
\end{eulerprompt}
\begin{eulerudf}
plot2d(none,r=1,title="Drag with the mouse, or press any key!");
start=0;
repeat
    {flag,m,time}=mousedrag();
    if flag==0 then return; endif;
    if flag==2 then
        hold on; mark(m[1],m[2]); hold off;
    endif;
end
endfunction
\end{eulerudf}
\begin{eulerprompt}
>dragtest // lihat hasilnya dan cobalah lakukan!
\end{eulerprompt}

\eulersubheading{Gaya pada Plot 2D}
\begin{eulercomment}
Secara bawaan, EMT menghitung tanda sumbu otomatis dan menambahkan label ke setiap tanda. Hal ini dapat diubah dengan parameter \texttt{grid}. Gaya bawaan sumbu dan labelnya dapat dimodifikasi. Selain itu, label dan judul dapat ditambahkan secara manual. Untuk mengatur ulang ke gaya bawaan, gunakan \texttt{reset()}.
\end{eulercomment}
\begin{eulerprompt}
>aspect();
>figure(3,4); ...
> figure(1); plot2d("x^3-x",grid=0); ... // tanpa grid, bingkai, atau sumbu
> figure(2); plot2d("x^3-x",grid=1); ... // sumbu-x-y
> figure(3); plot2d("x^3-x",grid=2); ... // tanda sumbu bawaan
> figure(4); plot2d("x^3-x",grid=3); ... // sumbu-x-y dengan label di dalam
> figure(5); plot2d("x^3-x",grid=4); ... // tanpa tanda sumbu, hanya label
> figure(6); plot2d("x^3-x",grid=5); ... // bawaan, tapi tanpa margin
> figure(7); plot2d("x^3-x",grid=6); ... // sumbu saja
> figure(8); plot2d("x^3-x",grid=7); ... // sumbu saja, tanda pada sumbu
> figure(9); plot2d("x^3-x",grid=8); ... // sumbu saja, tanda yang lebih halus pada sumbu
> figure(10); plot2d("x^3-x",grid=9); ... // bawaan, tanda kecil di dalam
> figure(11); plot2d("x^3-x",grid=10); ...// tanpa tanda sumbu, sumbu saja
> figure(0):
\end{eulerprompt}
\eulerimg{25}{images/Tugas 3-073.png}
\begin{eulercomment}
Parameter \texttt{<frame} mematikan bingkai, dan \texttt{framecolor=blue} mengatur bingkai menjadi warna biru.

Jika Anda menginginkan tanda sumbu (\emph{ticks}) Anda sendiri, Anda dapat menggunakan \texttt{grid=0}, dan menambahkan semuanya nanti.
\end{eulercomment}
\begin{eulerprompt}
>aspect(1.5); 
>plot2d("x^3-x",grid=0); // gambar plot
>frame; xgrid([-1,0,1]); ygrid(0): // tambahkan bingkai dan grid
\end{eulerprompt}
\eulerimg{30}{images/Tugas 3-074.png}
\begin{eulercomment}
Untuk judul plot dan label sumbu, lihat contoh berikut.
\end{eulercomment}
\begin{eulerprompt}
>plot2d("exp(x)",-1,1);
>textcolor(black); // atur warna teks menjadi hitam
>title(latex("y=e^x")); // judul di atas plot
>xlabel(latex("x")); // "x" untuk sumbu-x
>ylabel(latex("y"),>vertical); // "y" vertikal untuk sumbu-y
>label(latex("(0,1)"),0,1,color=blue): // beri label pada titik
\end{eulerprompt}
\eulerimg{30}{images/Tugas 3-075.png}
\begin{eulercomment}
Sumbu-sumbu dapat digambar secara terpisah dengan \texttt{xaxis()} dan \texttt{yaxis()}.
\end{eulercomment}
\begin{eulerprompt}
>plot2d("x^3-x",<grid,<frame);
>xaxis(0,xx=-2:1,style="->"); yaxis(0,yy=-5:5,style="->"):
\end{eulerprompt}
\eulerimg{30}{images/Tugas 3-076.png}
\begin{eulercomment}
Teks pada plot dapat diatur dengan \texttt{label()}. Pada contoh berikut, \texttt{"lc"} berarti \emph{lower center} (tengah bawah). Ini mengatur posisi label relatif terhadap koordinat plot.
\end{eulercomment}
\begin{eulerprompt}
>function f(x) &= x^3-x
\end{eulerprompt}
\begin{euleroutput}
        3
       x  - x
\end{euleroutput}
\begin{eulerprompt}
>plot2d(f,-1,1,>square);
>x0=fmin(f,0,1); // hitung titik minimum
>label("Rel. Min.",x0,f(x0),pos="lc"): // tambahkan label di sana
\end{eulerprompt}
\eulerimg{30}{images/Tugas 3-077.png}
\begin{eulercomment}
Terdapat juga kotak-kotak teks.
\end{eulercomment}
\begin{eulerprompt}
>plot2d(&f(x),-1,1,-2,2); // function
>plot2d(&diff(f(x),x),>add,style="--",color=red); // derivative
>labelbox(["f","f'"],["-","--"],[black,red]): // kotak label
\end{eulerprompt}
\eulerimg{30}{images/Tugas 3-078.png}
\begin{eulerprompt}
>plot2d(["exp(x)","1+x"],color=[black,blue],style=["-","-.-"]):
\end{eulerprompt}
\eulerimg{30}{images/Tugas 3-079.png}
\begin{eulerprompt}
>gridstyle("->",color=gray,textcolor=gray,framecolor=gray); ...
> plot2d("x^3-x",grid=1); ...
> settitle("y=x^3-x",color=black); ...
> label("x",2,0,pos="bc",color=gray); ...
> label("y",0,6,pos="cl",color=gray); ...
> reset():
\end{eulerprompt}
\eulerimg{25}{images/Tugas 3-080.png}
\begin{eulercomment}
Untuk kendali yang lebih jauh, sumbu-$x$ dan sumbu-$y$ dapat dibuat secara manual.

Perintah \texttt{fullwindow()} memperluas jendela plot karena kita tidak lagi membutuhkan tempat untuk label di luar jendela plot. Gunakan \texttt{shrinkwindow()} atau \texttt{reset()} untuk mengembalikannya ke pengaturan bawaan.
\end{eulercomment}
\begin{eulerprompt}
>fullwindow; ...
> gridstyle(color=darkgray,textcolor=darkgray); ...
> plot2d(["2^x","1","2^(-x)"],a=-2,b=2,c=0,d=4,<grid,color=4:6,<frame); ...
> xaxis(0,-2:1,style="->"); xaxis(0,2,"x",<axis); ...
> yaxis(0,4,"y",style="->"); ...
> yaxis(-2,1:4,>left); ...
> yaxis(2,2^(-2:2),style=".",<left); ...
> labelbox(["2^x","1","2^-x"],colors=4:6,x=0.8,y=0.2); ...
> reset:
\end{eulerprompt}
\eulerimg{25}{images/Tugas 3-081.png}
\begin{eulercomment}
Berikut adalah contoh lain, di mana string Unicode digunakan dan sumbu berada di luar area plot.
\end{eulercomment}
\begin{eulerprompt}
>aspect(1.5); 
>plot2d(["sin(x)","cos(x)"],0,2pi,color=[red,green],<grid,<frame); ...
> xaxis(-1.1,(0:2)*pi,xt=["0",u"&pi;",u"2&pi;"],style="-",>ticks,>zero); ...
> xgrid((0:0.5:2)*pi,<ticks); ...
> yaxis(-0.1*pi,-1:0.2:1,style="-",>zero,>grid); ...
> labelbox(["sin","cos"],colors=[red,green],x=0.5,y=0.2,>left); ...
> xlabel(u"&phi;"); ylabel(u"f(&phi;)"):
\end{eulerprompt}
\eulerimg{30}{images/Tugas 3-082.png}

\eulerheading{Plotting Data 2D}
\begin{eulercomment}
Jika $x$ dan $y$ adalah vektor data, data ini akan digunakan sebagai koordinat-$x$ dan koordinat-$y$ dari sebuah kurva. Pada kasus ini, \texttt{a}, \texttt{b}, \texttt{c}, dan \texttt{d}, atau sebuah radius \texttt{r} dapat ditentukan, atau jendela plot akan menyesuaikan secara otomatis dengan data tersebut. Secara alternatif, \texttt{>square} dapat diatur untuk mempertahankan rasio aspek persegi.

Menggambar plot dari sebuah ekspresi hanyalah singkatan untuk plot data. Untuk plot data, Anda membutuhkan satu atau lebih baris nilai-$x$, dan satu atau lebih baris nilai-$y$. Dari rentang dan nilai-$x$, fungsi \texttt{plot2d} akan menghitung data untuk digambar, secara bawaan menggunakan evaluasi adaptif pada fungsi tersebut. Untuk plot titik gunakan \texttt{>points}, dan untuk campuran garis dan titik gunakan \texttt{>addpoints}.

Namun Anda dapat memasukkan data secara langsung:
\begin{itemize}
    \item Gunakan vektor baris untuk $x$ dan $y$ pada satu fungsi.
    \item Matriks untuk $x$ dan $y$ digambar baris demi baris.
\end{itemize}

Berikut adalah contoh dengan satu baris untuk $x$ dan $y$.
\end{eulercomment}
\begin{eulerprompt}
>x=-10:0.1:10; y=exp(-x^2)*x; plot2d(x,y):
\end{eulerprompt}
\eulerimg{30}{images/Tugas 3-083.png}
\begin{eulercomment}
Data juga dapat digambar sebagai titik-titik. Gunakan \texttt{points=true} untuk hal ini. Plot ini berfungsi seperti poligon, namun hanya menggambar sudut-sudutnya.
\begin{itemize}
    \item \texttt{style="..."}: Pilih dari \texttt{"[]"}, \texttt{"<>"}, \texttt{"o"}, \texttt{"."}, \texttt{".."}, \texttt{"+"}, \texttt{"*"}, \texttt{"[]\#"}, \texttt{"<>\#"}, \texttt{"o\#"}, \texttt{"..\#"}, \texttt{"\#"}, \texttt{"|"}.
\end{itemize}

Untuk menggambar kumpulan titik, gunakan \texttt{>points}. Jika warnanya adalah vektor warna, setiap titik mendapatkan warna yang berbeda. Untuk matriks koordinat dan vektor kolom, warnanya berlaku pada baris-baris matriks tersebut.
Parameter \texttt{>addpoints} menambahkan titik-titik ke segmen garis untuk plot data.
\end{eulercomment}
\begin{eulerprompt}
>xdata=[1,1.5,2.5,3,4]; ydata=[3,3.1,2.8,2.9,2.7]; // data
>plot2d(xdata,ydata,a=0.5,b=4.5,c=2.5,d=3.5,style="."); // garis
>plot2d(xdata,ydata,>points,>add,style="o"): // tambahkan titik
\end{eulerprompt}
\eulerimg{30}{images/Tugas 3-084.png}
\begin{eulerprompt}
>p=polyfit(xdata,ydata,1); // dapatkan garis regresi
>plot2d("polyval(p,x)",>add,color=red): // tambahkan plot garis
\end{eulerprompt}
\eulerimg{30}{images/Tugas 3-085.png}

\eulerheading{Menggambar Daerah Yang Dibatasi Kurva}
\begin{eulercomment}
Plot data pada dasarnya adalah poligon. Kita juga dapat menggambar kurva atau kurva terisi (\emph{filled curves}).
\begin{itemize}
    \item \texttt{filled=true} mengisi plot tersebut.
    \item \texttt{style="..."}: Pilih dari \texttt{"\#"}, \texttt{"/"}, \texttt{"\textbackslash{}"}, \texttt{"\textbackslash{}/"}.
    \item \texttt{fillcolor}: Lihat di atas untuk warna yang tersedia.
\end{itemize}

Warna isian ditentukan oleh argumen \texttt{"fillcolor"}, dan \texttt{<outline} opsional mencegah penggambaran batas untuk semua gaya kecuali gaya bawaan.
\end{eulercomment}
\begin{eulerprompt}
>t=linspace(0,2pi,1000); // parameter untuk kurva
>x=sin(t)*exp(t/pi); y=cos(t)*exp(t/pi); // x(t) dan y(t)
>figure(1,2); aspect(16/9)
>figure(1); plot2d(x,y,r=10); // gambar kurva
>figure(2); plot2d(x,y,r=10,>filled,style="/",fillcolor=red); // isi kurva
>figure(0):
\end{eulerprompt}
\eulerimg{25}{images/Tugas 3-086.png}
\begin{eulercomment}
Pada contoh-contoh berikut kita menggambar elips terisi dan dua heksagon terisi menggunakan kurva tertutup dengan 6 titik dan gaya isian yang berbeda.
\end{eulercomment}
\begin{eulerprompt}
>x=linspace(0,2pi,1000); plot2d(sin(x),cos(x)*0.5,r=1,>filled,style="/"):
\end{eulerprompt}
\eulerimg{25}{images/Tugas 3-087.png}
\begin{eulerprompt}
>t=linspace(0,2pi,6); ...
>plot2d(cos(t),sin(t),>filled,style="/",fillcolor=red,r=1.2):
\end{eulerprompt}
\eulerimg{25}{images/Tugas 3-088.png}
\begin{eulerprompt}
>t=linspace(0,2pi,6); plot2d(cos(t),sin(t),>filled,style="#"):
\end{eulerprompt}
\eulerimg{25}{images/Tugas 3-089.png}
\begin{eulercomment}
Contoh lain adalah septagon (segi tujuh), yang kita buat dengan 7 titik pada lingkaran satuan.
\end{eulercomment}
\begin{eulerprompt}
>t=linspace(0,2pi,7); ...
> plot2d(cos(t),sin(t),r=1,>filled,style="/",fillcolor=red):
\end{eulerprompt}
\eulerimg{25}{images/Tugas 3-090.png}
\begin{eulercomment}
Berikut ini adalah himpunan nilai maksimal dari empat kondisi linear kurang dari atau sama dengan 3. Ini adalah $A[k] \cdot v \le 3$ untuk semua baris $A$. Untuk mendapatkan sudut yang bagus, kita menggunakan $n$ yang relatif besar.
\end{eulercomment}
\begin{eulerprompt}
>A=[2,1;1,2;-1,0;0,-1];
>function f(x,y) := max([x,y].A');
>plot2d("f",r=4,level=[0;3],color=green,n=111):
\end{eulerprompt}
\eulerimg{25}{images/Tugas 3-091.png}
\begin{eulercomment}
Poin utama dari bahasa matriks adalah ia memungkinkan kita menghasilkan tabel dari fungsi-fungsi secara mudah.
\end{eulercomment}
\begin{eulerprompt}
>t=linspace(0,2pi,1000); x=cos(3*t); y=sin(4*t);
\end{eulerprompt}
\begin{eulercomment}
Kita sekarang memiliki vektor nilai $x$ dan $y$. \texttt{plot2d()} dapat menggambar nilai-nilai ini sebagai kurva yang menghubungkan titik-titik tersebut. Plot ini dapat diisi. Pada kasus ini, hal tersebut menghasilkan gambaran yang bagus berkat aturan penggulungan (\emph{winding rule}) yang digunakan untuk isiannya.
\end{eulercomment}
\begin{eulerprompt}
>plot2d(x,y,<grid,<frame,>filled):
\end{eulerprompt}
\eulerimg{25}{images/Tugas 3-092.png}
\begin{eulercomment}
Vektor interval digambar terhadap nilai-$x$ sebagai area terisi di antara nilai batas bawah dan batas atas dari interval tersebut.

Hal ini dapat berguna untuk menggambar galat komputasi. Namun ini juga dapat digunakan untuk menggambar galat statistik.
\end{eulercomment}
\begin{eulerprompt}
>t=0:0.1:1; ...
> plot2d(t,interval(t-random(size(t)),t+random(size(t))),style="|"); ...
> plot2d(t,t,add=true):
\end{eulerprompt}
\eulerimg{25}{images/Tugas 3-093.png}
\begin{eulercomment}
Jika $x$ adalah vektor yang diurutkan, dan $y$ adalah vektor interval, maka \texttt{plot2d} akan menggambar rentang interval yang terisi pada bidang tersebut. Gaya isiannya sama dengan gaya poligon.
\end{eulercomment}
\begin{eulerprompt}
>t=-1:0.01:1; x=~t-0.01,t+0.01~; y=x^3-x;
>plot2d(t,y):
\end{eulerprompt}
\eulerimg{25}{images/Tugas 3-094.png}
\begin{eulercomment}
Dimungkinkan untuk mengisi area nilai untuk fungsi tertentu. Untuk ini, parameter \texttt{level} harus berupa matriks $2 \times n$. Baris pertama adalah batas bawah dan baris kedua berisi batas atas.
\end{eulercomment}
\begin{eulerprompt}
>expr := "2*x^2+x*y+3*y^4+y"; // definisikan ekspresi f(x,y)
>plot2d(expr,level=[0;1],style="-",color=blue): // 0 <= f(x,y) <= 1
\end{eulerprompt}
\eulerimg{25}{images/Tugas 3-095.png}
\begin{eulercomment}
Kita juga dapat mengisi rentang nilai seperti
\end{eulercomment}
\begin{eulerformula}
\[
-1 \le (x^2+y^2)^2-x^2+y^2 \le 0.
\]
\end{eulerformula}
\begin{eulerprompt}
>plot2d("(x^2+y^2)^2-x^2+y^2",r=1.2,level=[-1;0],style="/"):
\end{eulerprompt}
\eulerimg{25}{images/Tugas 3-097.png}
\begin{eulerprompt}
>plot2d("cos(x)","sin(x)^3",xmin=0,xmax=2pi,>filled,style="/"):
\end{eulerprompt}
\eulerimg{25}{images/Tugas 3-098.png}

\eulerheading{Grafik Fungsi Parametrik}
\begin{eulercomment}
Nilai-$x$ tidak perlu diurutkan. $(x,y)$ secara sederhana mendeskripsikan sebuah kurva. Jika $x$ diurutkan, maka kurva tersebut adalah grafik dari sebuah fungsi.

Pada contoh berikut, kita menggambar sebuah spiral
\end{eulercomment}
\begin{eulerformula}
\[
\gamma(t) = t \cdot (\cos(2\pi t),\sin(2\pi t))
\]
\end{eulerformula}
\begin{eulercomment}
Kita perlu menggunakan titik yang sangat banyak agar terlihat mulus, atau menggunakan fungsi \texttt{adaptive()} untuk mengevaluasi ekspresinya (lihat fungsi \texttt{adaptive()} untuk detail lebih lanjut).
\end{eulercomment}
\begin{eulerprompt}
>t=linspace(0,1,1000); ...
>plot2d(t*cos(2*pi*t),t*sin(2*pi*t),r=1):
\end{eulerprompt}
\eulerimg{25}{images/Tugas 3-100.png}
\begin{eulercomment}
Sebagai alternatif, sangat memungkinkan untuk menggunakan dua ekspresi pada kurva. Hal berikut ini menggambar kurva yang sama seperti di atas.
\end{eulercomment}
\begin{eulerprompt}
>plot2d("x*cos(2*pi*x)","x*sin(2*pi*x)",xmin=0,xmax=1,r=1):
\end{eulerprompt}
\eulerimg{25}{images/Tugas 3-101.png}
\begin{eulerprompt}
>t=linspace(0,1,1000); r=exp(-t); x=r*cos(2pi*t); y=r*sin(2pi*t);
>plot2d(x,y,r=1):
\end{eulerprompt}
\eulerimg{25}{images/Tugas 3-102.png}
\begin{eulercomment}
Pada contoh selanjutnya, kita menggambar kurva
\end{eulercomment}
\begin{eulerformula}
\[
\gamma(t) = (r(t) \cos(t), r(t) \sin(t))
\]
\end{eulerformula}
\begin{eulercomment}
dengan
\end{eulercomment}
\begin{eulerformula}
\[
r(t) = 1 + \dfrac{\sin(3t)}{2}.
\]
\end{eulerformula}
\begin{eulerprompt}
>t=linspace(0,2pi,1000); r=1+sin(3*t)/2; x=r*cos(t); y=r*sin(t); ...
>plot2d(x,y,>filled,fillcolor=red,style="/",r=1.5):
\end{eulerprompt}
\eulerimg{25}{images/Tugas 3-105.png}

\eulerheading{Menggambar Grafik Bilangan Kompleks}
\begin{eulercomment}
Larik bilangan kompleks juga dapat digambar. Titik-titik grid nantinya akan saling terhubung. Jika sejumlah garis grid ditentukan (atau vektor $1 \times 2$ untuk garis grid) pada argumen \texttt{cgrid}, hanya garis grid tersebut yang akan terlihat.

Matriks bilangan kompleks akan otomatis tergambar sebagai grid di bidang kompleks.

Pada contoh berikut, kita memplot bayangan dari lingkaran satuan di bawah fungsi eksponensial. Parameter \texttt{cgrid} menyembunyikan beberapa kurva grid.
\end{eulercomment}
\begin{eulerprompt}
>aspect(); r=linspace(0,1,50); a=linspace(0,2pi,80)'; z=r*exp(I*a);...
>plot2d(z,a=-1.25,b=1.25,c=-1.25,d=1.25,cgrid=10):
\end{eulerprompt}
\eulerimg{25}{images/Tugas 3-106.png}
\begin{eulerprompt}
>aspect(1.25); r=linspace(0,1,50); a=linspace(0,2pi,200)'; z=r*exp(I*a);
>plot2d(exp(z),cgrid=[40,10]):
\end{eulerprompt}
\eulerimg{34}{images/Tugas 3-107.png}
\begin{eulerprompt}
>r=linspace(0,1,10); a=linspace(0,2pi,40)'; z=r*exp(I*a);
>plot2d(exp(z),>points,>add):
\end{eulerprompt}
\eulerimg{34}{images/Tugas 3-108.png}
\begin{eulercomment}
Vektor bilangan kompleks secara otomatis akan digambar sebagai kurva di bidang kompleks dengan bagian real dan bagian imajinernya.

Pada contoh ini, kita menggambar kurva kompleks berdasarkan lingkaran satuan
\end{eulercomment}
\begin{eulerformula}
\[
\gamma(t) = e^{it}
\]
\end{eulerformula}
\begin{eulerprompt}
>t=linspace(0,2pi,1000); ...
>plot2d(exp(I*t)+exp(4*I*t),r=2):
\end{eulerprompt}
\eulerimg{34}{images/Tugas 3-110.png}

\eulerheading{Plot Statistik}
\begin{eulercomment}
Ada banyak fungsi yang dikhususkan untuk plot statistik. Salah satu yang paling sering digunakan adalah plot kolom (\emph{column plot}).

Jumlah kumulatif dari nilai terdistribusi normal $0$-$1$ akan menghasilkan jalan acak (\emph{random walk}).
\end{eulercomment}
\begin{eulerprompt}
>plot2d(cumsum(randnormal(1,1000))):
\end{eulerprompt}
\eulerimg{34}{images/Tugas 3-111.png}
\begin{eulercomment}
Penggunaan dua baris akan menunjukkan pergerakan dalam dua dimensi.
\end{eulercomment}
\begin{eulerprompt}
>X=cumsum(randnormal(2,1000)); plot2d(X[1],X[2]):
\end{eulerprompt}
\eulerimg{34}{images/Tugas 3-112.png}
\begin{eulerprompt}
>columnsplot(cumsum(random(10)),style="/",color=blue):
\end{eulerprompt}
\eulerimg{34}{images/Tugas 3-113.png}
\begin{eulercomment}
Hal ini juga dapat menampilkan string sebagai label.
\end{eulercomment}
\begin{eulerprompt}
>months=["Jan","Feb","Mar","Apr","May","Jun", ...
>  "Jul","Aug","Sep","Oct","Nov","Dec"];
>values=[10,12,12,18,22,28,30,26,22,18,12,8];
>columnsplot(values,lab=months,color=red,style="-");
>title("Temperature"):
\end{eulerprompt}
\eulerimg{34}{images/Tugas 3-114.png}
\begin{eulerprompt}
>k=0:10;
>plot2d(k,bin(10,k),>bar):
\end{eulerprompt}
\eulerimg{34}{images/Tugas 3-115.png}
\begin{eulerprompt}
>plot2d(k,bin(10,k)); plot2d(k,bin(10,k),>points,>add):
\end{eulerprompt}
\eulerimg{34}{images/Tugas 3-116.png}
\begin{eulerprompt}
>plot2d(normal(1000),normal(1000),>points,grid=6,style=".."):
\end{eulerprompt}
\eulerimg{34}{images/Tugas 3-117.png}
\begin{eulerprompt}
>plot2d(normal(1,1000),>distribution,style="O"):
\end{eulerprompt}
\eulerimg{34}{images/Tugas 3-118.png}
\begin{eulerprompt}
>plot2d("qnormal",0,5;2.5,0.5,>filled):
\end{eulerprompt}
\eulerimg{34}{images/Tugas 3-119.png}
\begin{eulercomment}
Untuk menggambar plot distribusi statistik eksperimental, Anda dapat menggunakan \texttt{distribution=n} dengan \texttt{plot2d}.
\end{eulercomment}
\begin{eulerprompt}
>w=randexponential(1,1000); // exponential distribution
>plot2d(w,>distribution): // or distribution=n with n intervals
\end{eulerprompt}
\eulerimg{34}{images/Tugas 3-120.png}
\begin{eulercomment}
Atau Anda dapat menghitung distribusi dari data dan memplot hasilnya dengan \texttt{>bar} di \texttt{plot2d}, atau dengan plot kolom.
\end{eulercomment}
\begin{eulerprompt}
>w=normal(1000); // 0-1-normal distribution
>{x,y}=histo(w,10,v=[-6,-4,-2,-1,0,1,2,4,6]); // batas interval v
>plot2d(x,y,>bar):
\end{eulerprompt}
\eulerimg{34}{images/Tugas 3-121.png}
\begin{eulercomment}
Fungsi \texttt{statplot()} mengatur gaya dengan sebuah string sederhana.
\end{eulercomment}
\begin{eulerprompt}
>statplot(1:10,cumsum(random(10)),"b"):
\end{eulerprompt}
\eulerimg{34}{images/Tugas 3-122.png}
\begin{eulerprompt}
>n=10; i=0:n; ...
>plot2d(i,bin(n,i)/2^n,a=0,b=10,c=0,d=0.3); ...
>plot2d(i,bin(n,i)/2^n,points=true,style="ow",add=true,color=blue):
\end{eulerprompt}
\eulerimg{34}{images/Tugas 3-123.png}
\begin{eulercomment}
Selain itu, data dapat digambar sebagai plot batang. Dalam kasus ini, $x$ harus diurutkan dan satu elemen lebih panjang dari $y$. Batang tersebut akan membentang dari \texttt{x[i]} ke \texttt{x[i+1]} dengan nilai \texttt{y[i]}. Jika $x$ memiliki ukuran yang sama dengan $y$, ia akan diperpanjang satu elemen menggunakan jarak (\emph{spacing}) terakhir.

Gaya isian dapat digunakan sama seperti yang di atas.
\end{eulercomment}
\begin{eulerprompt}
>n=10; k=bin(n,0:n); ...
>plot2d(-0.5:n+0.5,k,bar=true,fillcolor=lightgray):
\end{eulerprompt}
\eulerimg{34}{images/Tugas 3-124.png}
\begin{eulercomment}
Data untuk plot batang (\texttt{bar=1}) dan histogram (\texttt{histogram=1}) dapat diberikan secara eksplisit pada \texttt{xv} dan \texttt{yv}, atau dapat dihitung dari distribusi empiris pada \texttt{xv} dengan \texttt{>distribution} (atau \texttt{distribution=n}). Histogram dari nilai-nilai \texttt{xv} akan dihitung secara otomatis dengan \texttt{>histogram}. Jika \texttt{>even} ditentukan, nilai-nilai \texttt{xv} akan dihitung dalam interval bilangan bulat.
\end{eulercomment}
\begin{eulerprompt}
>plot2d(normal(10000),distribution=50):
\end{eulerprompt}
\eulerimg{34}{images/Tugas 3-125.png}
\begin{eulerprompt}
>k=0:10; m=bin(10,k); x=(0:11)-0.5; plot2d(x,m,>bar):
\end{eulerprompt}
\eulerimg{34}{images/Tugas 3-126.png}
\begin{eulerprompt}
>columnsplot(m,k):
\end{eulerprompt}
\eulerimg{34}{images/Tugas 3-127.png}
\begin{eulerprompt}
>plot2d(random(600)*6,histogram=6):
\end{eulerprompt}
\eulerimg{34}{images/Tugas 3-128.png}
\begin{eulercomment}
Untuk distribusi, terdapat parameter \texttt{distribution=n}, yang menghitung nilai-nilai secara otomatis dan mencetak distribusi relatif dengan $n$ sub-interval.
\end{eulercomment}
\begin{eulerprompt}
>plot2d(normal(1,1000),distribution=10,style="\/"):
\end{eulerprompt}
\eulerimg{34}{images/Tugas 3-129.png}
\begin{eulercomment}
Dengan parameter \texttt{even=true}, ini akan menggunakan interval bilangan bulat.
\end{eulercomment}
\begin{eulerprompt}
>plot2d(intrandom(1,1000,10),distribution=10,even=true):
\end{eulerprompt}
\eulerimg{34}{images/Tugas 3-130.png}
\begin{eulercomment}
Perhatikan bahwa ada banyak plot statistik yang mungkin berguna. Silakan lihat tutorial mengenai statistik.
\end{eulercomment}
\begin{eulerprompt}
>columnsplot(getmultiplicities(1:6,intrandom(1,6000,6))):
\end{eulerprompt}
\eulerimg{34}{images/Tugas 3-131.png}
\begin{eulerprompt}
>plot2d(normal(1,1000),>distribution); ...
>  plot2d("qnormal(x)",color=red,thickness=2,>add):
\end{eulerprompt}
\eulerimg{34}{images/Tugas 3-132.png}
\begin{eulercomment}
Ada banyak juga plot khusus untuk statistik. Boxplot menunjukkan kuartil dari distribusi ini dan banyaknya pencilan (\emph{outlier}). Berdasarkan definisi, pencilan pada boxplot adalah data yang melebihi $1{,}5$ kali rentang $50\%$ bagian tengah dari plot tersebut.
\end{eulercomment}
\begin{eulerprompt}
>M=normal(5,1000); boxplot(quartiles(M)):
\end{eulerprompt}
\eulerimg{34}{images/Tugas 3-133.png}

\eulerheading{Fungsi Implisit}
\begin{eulercomment}
Plot implisit menunjukkan garis-garis level (\emph{level lines}) yang menyelesaikan $f(x,y)=\text{level}$, di mana \texttt{"level"} bisa berupa satu nilai atau vektor dari nilai. Jika \texttt{level="auto"}, akan ada \texttt{nc} garis level, yang akan tersebar secara merata di antara minimum dan maksimum fungsi tersebut. Warna yang lebih gelap atau lebih terang dapat ditambahkan dengan \texttt{>hue} untuk menunjukkan nilai fungsi tersebut. Untuk fungsi implisit, \texttt{xv} harus berupa fungsi atau ekspresi dari parameter $x$ dan $y$, atau, sebagai alternatif, \texttt{xv} bisa berupa matriks nilai.

Euler dapat menandai garis level
\end{eulercomment}
\begin{eulerformula}
\[
f(x,y) = c
\]
\end{eulerformula}
\begin{eulercomment}
dari sembarang fungsi.

Untuk menggambar himpunan $f(x,y)=c$ dari satu atau lebih konstanta $c$, Anda dapat menggunakan \texttt{plot2d()} dengan plot implisitnya di bidang. Parameter untuk $c$ adalah \texttt{level=c}, di mana $c$ dapat berupa vektor garis level. Selain itu, skema warna dapat digambar di latar belakang untuk menunjukkan nilai dari fungsi tersebut untuk tiap titik pada plot. Parameter \texttt{"n"} menentukan tingkat kehalusan plot tersebut.
\end{eulercomment}
\begin{eulerprompt}
>aspect(1.5); 
>plot2d("x^2+y^2-x*y-x",r=1.5,level=0,contourcolor=red):
\end{eulerprompt}
\eulerimg{30}{images/Tugas 3-135.png}
\begin{eulerprompt}
>expr := "2*x^2+x*y+3*y^4+y"; // definisikan ekspresi f(x,y)
>plot2d(expr,level=0): // Solusi dari f(x,y)=0
\end{eulerprompt}
\eulerimg{30}{images/Tugas 3-136.png}
\begin{eulerprompt}
>plot2d(expr,level=0:0.5:20,>hue,contourcolor=white,n=200): // bagus
\end{eulerprompt}
\eulerimg{30}{images/Tugas 3-137.png}
\begin{eulerprompt}
>plot2d(expr,level=0:0.5:20,>hue,>spectral,n=200,grid=4): // lebih bagus
\end{eulerprompt}
\eulerimg{30}{images/Tugas 3-138.png}
\begin{eulercomment}
Hal ini berfungsi untuk plot data juga. Namun Anda harus menentukan rentangnya untuk label sumbunya.
\end{eulercomment}
\begin{eulerprompt}
>x=-2:0.05:1; y=x'; z=expr(x,y);
>plot2d(z,level=0,a=-1,b=2,c=-2,d=1,>hue):
\end{eulerprompt}
\eulerimg{30}{images/Tugas 3-139.png}
\begin{eulerprompt}
>plot2d("x^3-y^2",>contour,>hue,>spectral):
\end{eulerprompt}
\eulerimg{30}{images/Tugas 3-140.png}
\begin{eulerprompt}
>plot2d("x^3-y^2",level=0,contourwidth=3,>add,contourcolor=red):
\end{eulerprompt}
\eulerimg{30}{images/Tugas 3-141.png}
\begin{eulerprompt}
>z=z+normal(size(z))*0.2;
>plot2d(z,level=0.5,a=-1,b=2,c=-2,d=1):
\end{eulerprompt}
\eulerimg{30}{images/Tugas 3-142.png}
\begin{eulerprompt}
>plot2d(expr,level=[0:0.2:5;0.05:0.2:5.05],color=lightgray):
\end{eulerprompt}
\eulerimg{30}{images/Tugas 3-143.png}
\begin{eulerprompt}
>plot2d("x^2+y^3+x*y",level=1,r=4,n=100):
\end{eulerprompt}
\eulerimg{30}{images/Tugas 3-144.png}
\begin{eulerprompt}
>plot2d("x^2+2*y^2-x*y",level=0:0.1:10,n=100,contourcolor=white,>hue):
\end{eulerprompt}
\eulerimg{30}{images/Tugas 3-145.png}
\begin{eulercomment}
Sangat memungkinkan juga untuk mengisi himpunan
\end{eulercomment}
\begin{eulerformula}
\[
a \le f(x,y) \le b
\]
\end{eulerformula}
\begin{eulercomment}
dengan rentang level.

Dimungkinkan untuk mengisi area nilai dari fungsi tertentu. Untuk ini, parameter \texttt{level} harus berupa matriks $2 \times n$. Baris pertama adalah batas bawah dan baris kedua berisi batas atas.
\end{eulercomment}
\begin{eulerprompt}
>plot2d(expr,level=[0;1],style="-",color=blue): // 0 <= f(x,y) <= 1
\end{eulerprompt}
\eulerimg{30}{images/Tugas 3-147.png}
\begin{eulercomment}
Plot implisit juga dapat menunjukkan rentang dari beberapa level. Maka parameter \texttt{level} harus berupa matriks $2 \times n$ untuk interval level, di mana baris pertamanya berisi awal dan baris keduanya memuat akhir dari tiap interval. Sebagai alternatif, vektor baris sederhana dapat digunakan untuk \texttt{level}, dan parameter \texttt{dl} memperluas nilai level tersebut menjadi interval.
\end{eulercomment}
\begin{eulerprompt}
>plot2d("x^4+y^4",r=1.5,level=[0;1],color=blue,style="/"):
\end{eulerprompt}
\eulerimg{30}{images/Tugas 3-148.png}
\begin{eulerprompt}
>plot2d("x^2+y^3+x*y",level=[0,2,4;1,3,5],style="/",r=2,n=100):
\end{eulerprompt}
\eulerimg{30}{images/Tugas 3-149.png}
\begin{eulerprompt}
>plot2d("x^2+y^3+x*y",level=-10:20,r=2,style="-",dl=0.1,n=100):
\end{eulerprompt}
\eulerimg{30}{images/Tugas 3-150.png}
\begin{eulerprompt}
>plot2d("sin(x)*cos(y)",r=pi,>hue,>levels,n=100):
\end{eulerprompt}
\eulerimg{30}{images/Tugas 3-151.png}
\begin{eulercomment}
Sangat mungkin juga untuk menandai area
\end{eulercomment}
\begin{eulerformula}
\[
a \le f(x,y) \le b.
\]
\end{eulerformula}
\begin{eulercomment}
Hal ini dilakukan dengan menambahkan sebuah level yang memiliki dua baris.
\end{eulercomment}
\begin{eulerprompt}
>plot2d("(x^2+y^2-1)^3-x^2*y^3",r=1.3, ...
>  style="#",color=red,<outline, ...
>  level=[-2;0],n=100):
\end{eulerprompt}
\eulerimg{30}{images/Tugas 3-153.png}
\begin{eulercomment}
Dimungkinkan untuk menentukan sebuah level yang spesifik. Misal, kita dapat memplot penyelesaian dari suatu persamaan seperti
\end{eulercomment}
\begin{eulerformula}
\[
x^3-xy+x^2y^2=6
\]
\end{eulerformula}
\begin{eulerprompt}
>plot2d("x^3-x*y+x^2*y^2",r=6,level=1,n=100):
\end{eulerprompt}
\eulerimg{30}{images/Tugas 3-155.png}
\begin{eulerprompt}
>function starplot1 (v, style="/", color=green, lab=none) ...
\end{eulerprompt}
\begin{eulerudf}
if !holding() then clg; endif;
w=window(); window(0,0,1024,1024);
h=holding(1);
r=max(abs(v))*1.2;
setplot(-r,r,-r,r);
n=cols(v); t=linspace(0,2pi,n);
v=v|v[1]; c=v*cos(t); s=v*sin(t);
cl=barcolor(color); st=barstyle(style);
loop 1 to n
    polygon([0,c[#],c[#+1]],[0,s[#],s[#+1]],1);
    if lab!=none then
        rlab=v[#]+r*0.1;
        {col,row}=toscreen(cos(t[#])*rlab,sin(t[#])*rlab);
        ctext(""+lab[#],col,row-textheight()/2);
    endif;
end;
barcolor(cl); barstyle(st);
holding(h);
window(w);
endfunction
\end{eulerudf}
\begin{eulercomment}
Tidak ada grid atau tanda sumbu di sini. Selain itu, kita menggunakan jendela penuh untuk plot ini.

Kita memanggil \texttt{reset} sebelum kita menguji plot ini untuk mengembalikan pengaturan bawaan grafisnya. Ini tidak diperlukan jika Anda yakin bahwa plot Anda sudah berfungsi dengan baik.
\end{eulercomment}
\begin{eulerprompt}
>reset; starplot1(normal(1,10)+5,color=red,lab=1:10):
\end{eulerprompt}
\eulerimg{25}{images/Tugas 3-156.png}
\begin{eulercomment}
Terkadang, Anda mungkin ingin memplot sesuatu yang tidak dapat dilakukan oleh \texttt{plot2d} secara persis, namun hampir bisa.

Pada fungsi berikut, kita membuat plot impuls logaritmik. \texttt{plot2d} bisa melakukan plot logaritmik, namun tidak untuk batang impuls.
\end{eulercomment}
\begin{eulerprompt}
>function logimpulseplot1 (x,y) ...
\end{eulerprompt}
\begin{eulerudf}
{x0,y0}=makeimpulse(x,log(y)/log(10));
plot2d(x0,y0,>bar,grid=0);
h=holding(1);
frame();
xgrid(ticks(x));
p=plot();
for i=-10 to 10;
    if i<=p[4] and i>=p[3] then
        ygrid(i,yt="10^"+i);
    endif;
end;
holding(h);
endfunction
\end{eulerudf}
\begin{eulercomment}
Mari kita mengujinya dengan nilai-nilai yang terdistribusi secara eksponensial.
\end{eulercomment}
\begin{eulerprompt}
>aspect(1.5); x=1:10; y=-log(random(size(x)))*200; ...
>logimpulseplot1(x,y):
\end{eulerprompt}
\eulerimg{30}{images/Tugas 3-157.png}
\begin{eulercomment}
Mari kita buat animasi pada kurva 2D menggunakan plot langsung (\emph{direct plots}). Perintah \texttt{plot(x,y)} secara sederhana akan memplot sebuah kurva ke dalam jendela plot. \texttt{setplot(a,b,c,d)} mengatur jendela ini.

Fungsi \texttt{wait(0)} memaksa plot untuk muncul pada jendela grafis. Jika tidak, penggambaran ulang terjadi dalam interval waktu yang lebih lambat/jarang.
\end{eulercomment}
\begin{eulerprompt}
>function animliss (n,m) ...
\end{eulerprompt}
\begin{eulerudf}
t=linspace(0,2pi,500);
f=0;
c=framecolor(0);
l=linewidth(2);
setplot(-1,1,-1,1);
repeat
    clg;
    plot(sin(n*t),cos(m*t+f));
    wait(0);
    if testkey() then break; endif;
    f=f+0.02;
end;
framecolor(c);
linewidth(l);
endfunction
\end{eulerudf}
\begin{eulercomment}
Tekan tombol mana saja untuk menghentikan animasi ini.
\end{eulercomment}
\begin{eulerprompt}
>animliss(2,3); // lihat hasilnya, jika sudah puas, tekan ENTER
\end{eulerprompt}

\eulerheading{Plot Logaritmik}
\begin{eulercomment}
EMT menggunakan parameter \texttt{"logplot"} untuk skala logaritmik.
Plot logaritmik dapat digambar baik menggunakan skala logaritmik di sumbu-$y$ dengan \texttt{logplot=1}, atau menggunakan skala logaritmik di sumbu-$x$ dan sumbu-$y$ dengan \texttt{logplot=2}, atau di sumbu-$x$ dengan \texttt{logplot=3}.
\end{eulercomment}
\begin{eulerttcomment}
- logplot=1: y-logaritmik
- logplot=2: x-y-logaritmik
- logplot=3: x-logaritmik
\end{eulerttcomment}
\begin{eulerprompt}
>plot2d("exp(x^3-x)*x^2",1,5,logplot=1):
\end{eulerprompt}
\eulerimg{30}{images/Tugas 3-158.png}
\begin{eulerprompt}
>plot2d("exp(x+sin(x))",0,100,logplot=1):
\end{eulerprompt}
\eulerimg{30}{images/Tugas 3-159.png}
\begin{eulerprompt}
>plot2d("exp(x+sin(x))",10,100,logplot=2):
\end{eulerprompt}
\eulerimg{30}{images/Tugas 3-160.png}
\begin{eulerprompt}
>plot2d("gamma(x)",1,10,logplot=1):
\end{eulerprompt}
\eulerimg{30}{images/Tugas 3-161.png}
\begin{eulerprompt}
>plot2d("log(x*(2+sin(x/100)))",10,1000,logplot=3):
\end{eulerprompt}
\eulerimg{30}{images/Tugas 3-162.png}
\begin{eulercomment}
Ini juga dapat berfungsi dengan plot data.
\end{eulercomment}
\begin{eulerprompt}
>x=10^(1:20); y=x^2-x;
>plot2d(x,y,logplot=2):
\end{eulerprompt}
\eulerimg{30}{images/Tugas 3-163.png}

\eulersubheading{Grafik Fungsi Satu Variabel dalam Bentuk Ekspresi Langsung}
\begin{eulerformula}
\[
y=x^2-4x+3
\]
\end{eulerformula}
\begin{eulerprompt}
>plot2d("x^2-4x+3",-10,10):
\end{eulerprompt}
\eulerimg{30}{images/Tugas 3-165.png}
\begin{eulerprompt}
>plot2d("cos(4*x)",-10,10):
\end{eulerprompt}
\eulerimg{30}{images/Tugas 3-166.png}
\begin{eulerprompt}
>plot2d("(cos((x^2)/4))^2",-5,5):
\end{eulerprompt}
\eulerimg{30}{images/Tugas 3-167.png}
\begin{eulerprompt}
>plot2d("sin(x^2+1)",-5,5):
\end{eulerprompt}
\eulerimg{30}{images/Tugas 3-168.png}

\eulersubheading{Grafik Fungsi Satu Variabel yang Rumusnya Disimpan dalam Variabel Ekspresi}
\begin{eulerformula}
\[
f(x)=x^2+1
\]
\end{eulerformula}
\begin{eulerprompt}
>reset; fx:="x^2+1";
>plot2d(fx,-5,5):
\end{eulerprompt}
\eulerimg{25}{images/Tugas 3-170.png}
\begin{eulerformula}
\[
g(x)=x^2+x-5
\]
\end{eulerformula}
\begin{eulerprompt}
>function g(x)&=x^2+x-5;
>plot2d("g",-5,5):
\end{eulerprompt}
\eulerimg{25}{images/Tugas 3-172.png}
\begin{eulerformula}
\[
h(x)=4x-x^3
\]
\end{eulerformula}
\begin{eulerprompt}
>hx="4*x-x^3";
>plot2d("h",-5,5):
\end{eulerprompt}
\eulerimg{25}{images/Tugas 3-174.png}

\eulersubheading{Grafik Fungsi Satu Variabel yang Fungsinya Didefinisikan sebagai Fungsi Numerik}
\begin{eulerformula}
\[
f(x)=x^2-4x+4
\]
\end{eulerformula}
\begin{eulerprompt}
>function f(x):=x^2-4*x+4
>plot2d("f",-2,6):
\end{eulerprompt}
\eulerimg{25}{images/Tugas 3-176.png}
\begin{eulerformula}
\[
f(x)=\sin(x)\cos(x)
\]
\end{eulerformula}
\begin{eulerprompt}
>function f(x):=sin(x)*cos(x)
>plot2d("f",-pi,pi):
\end{eulerprompt}
\eulerimg{25}{images/Tugas 3-178.png}

\eulersubheading{Grafik Fungsi Satu Variabel yang Fungsinya Didefinisikan sebagai Fungsi Simbolik}
\begin{eulerformula}
\[
f(x)=x^3-3x
\]
\end{eulerformula}
\begin{eulerprompt}
>function f(x)&=x^3-3*x;
>plot2d("f",-3,3):
\end{eulerprompt}
\eulerimg{25}{images/Tugas 3-180.png}
\begin{eulerformula}
\[
\frac{dy}{dx}, \quad y=\sin(x)e^x
\]
\end{eulerformula}
\begin{eulerprompt}
>plot2d(&diff(sin(x)*exp(x),x),-pi,pi):
\end{eulerprompt}
\eulerimg{25}{images/Tugas 3-182.png}

\eulersubheading{Grafik Beberapa Kurva Sekaligus}
\begin{eulerformula}
\[
y=x^2, \quad y=x^3, \quad y=x^4
\]
\end{eulerformula}
\begin{eulerprompt}
>plot2d(["x^2","x^3","x^4"],-2,2):
\end{eulerprompt}
\eulerimg{25}{images/Tugas 3-184.png}
\begin{eulerformula}
\[
y=\sin(x), \quad y=\cos(x)
\]
\end{eulerformula}
\begin{eulerprompt}
>x=linspace(-2,2,100);
>plot2d(x,[sin(x);cos(x)]):
\end{eulerprompt}
\eulerimg{25}{images/Tugas 3-186.png}

\eulersubheading{Grafik Beberapa Kurva dalam Satu Bidang Koordinat}
\begin{eulerformula}
\[
y=x^2, \quad y=x^3
\]
\end{eulerformula}
\begin{eulerprompt}
>plot2d("x^2",-2,2);
>plot2d("x^3",-2,2,>add):
\end{eulerprompt}
\eulerimg{25}{images/Tugas 3-188.png}
\begin{eulerformula}
\[
\text{Sebaran titik (scatter) dan kurva } y=x^2
\]
\end{eulerformula}
\begin{eulerprompt}
>plot2d(1:10,(1:10)^2+random(1,10),>points);
>plot2d("x^2",1,10,>add):
\end{eulerprompt}
\eulerimg{25}{images/Tugas 3-190.png}

\eulersubheading{Label Sumbu Koordinat, Label Kurva, dan Keterangan Kurva (Legend)}
\begin{eulerformula}
\[
y=\sin(x)
\]
\end{eulerformula}
\begin{eulercomment}
dengan keterangan sumbu-$x$ sebagai waktu, sumbu-$y$ sebagai amplitudo, dan judul, yaitu gelombang sinus.
\end{eulercomment}
\begin{eulerprompt}
>plot2d("sin(x)",0,2*pi,xl="waktu(t)",yl="amplitudo",title="gelombang sinus"):
\end{eulerprompt}
\eulerimg{25}{images/Tugas 3-192.png}
\begin{eulerformula}
\[
y=x^2
\]
\end{eulerformula}
\begin{eulercomment}
dengan label kurva di titik tertentu.
\end{eulercomment}
\begin{eulerprompt}
>plot2d("x^2",-2,2);
>label("titik minimum",0,0):
\end{eulerprompt}
\eulerimg{25}{images/Tugas 3-194.png}

\eulersubheading{Mengatur Ukuran Grafik, Format (Style), dan Warna Kurva}
\begin{eulerformula}
\[
y=\sin(x)
\]
\end{eulerformula}
\begin{eulercomment}
dengan warna merah dan ketebalan 3 satuan.
\end{eulercomment}
\begin{eulerprompt}
>plot2d("sin(x)",color=red,thickness=3):
\end{eulerprompt}
\eulerimg{25}{images/Tugas 3-196.png}
\begin{eulerformula}
\[
y=x^2
\]
\end{eulerformula}
\begin{eulercomment}
dengan warna biru dan garis putus-putus.
\end{eulercomment}
\begin{eulerprompt}
>plot2d("x^2",color=blue,style="--"):
\end{eulerprompt}
\eulerimg{25}{images/Tugas 3-198.png}

\eulersubheading{Grafik Sekumpulan Kurva dengan Satu Perintah plot2d}
\begin{eulerformula}
\[
y=ax^2, \quad a=[1,5]
\]
\end{eulerformula}
\begin{eulerprompt}
>function f(x,a):=a*x^2;
>plot2d("f(x,1)",-2,2);
>for a=2:5;plot2d("f(x,"+a+")",-2,2,>add); end:
\end{eulerprompt}
\eulerimg{25}{images/Tugas 3-200.png}
\begin{eulerprompt}
>plot2d("sin(x)",0,2*pi);
>for c=0:.5:2; plot2d("sin(x-"+c+")",0,2*pi,>add); end:
\end{eulerprompt}
\eulerimg{25}{images/Tugas 3-201.png}

\eulersubheading{Gambar Kurva Interaktif}
\begin{eulerformula}
\[
y=\sin\left(\frac{1}{x}\right)
\]
\end{eulerformula}
\begin{eulerprompt}
>plot2d("sin(1/x)",.01,1):
\end{eulerprompt}
\eulerimg{25}{images/Tugas 3-203.png}
\begin{eulerprompt}
>zoom():
\end{eulerprompt}
\eulerimg{25}{images/Tugas 3-204.png}
\begin{eulerprompt}
>mouse():
\end{eulerprompt}
\eulerimg{25}{images/Tugas 3-205.png}

\eulersubheading{Grafik Kurva Fungsi Parametrik}
\begin{eulerformula}
\[
x(t)=\cos(t), \quad y(t)=\sin(t)
\]
\end{eulerformula}
\begin{eulerprompt}
>t=linspace(0,2*pi,100);
>plot2d(cos(t),sin(t)):
\end{eulerprompt}
\eulerimg{25}{images/Tugas 3-207.png}
\begin{eulerprompt}
>t=linspace(0,2*pi,500);
>plot2d(sin(3*t),cos(4*t)):
\end{eulerprompt}
\eulerimg{25}{images/Tugas 3-208.png}

\eulersubheading{Grafik Kurva Fungsi Implisit}
\begin{eulerformula}
\[
F(x,y)=x^2+y^2-4=0
\]
\end{eulerformula}
\begin{eulerprompt}
>plot2d("x^2+y^2-4",r=3,level=0):
\end{eulerprompt}
\eulerimg{25}{images/Tugas 3-210.png}
\begin{eulerformula}
\[
F(x,y)=x^3+y^3-3xy=0
\]
\end{eulerformula}
\begin{eulerprompt}
>plot2d("x^3+y^3-3*x*y",r=3,level=0):
\end{eulerprompt}
\eulerimg{25}{images/Tugas 3-212.png}

\eulersubheading{Grafik Kurva Fungsi Kompleks}
\begin{eulerformula}
\[
z=e^{it}
\]
\end{eulerformula}
\begin{eulerprompt}
>t=linspace(0,2*pi,100);
>z=exp(%i*t);
>plot2d(z):
\end{eulerprompt}
\eulerimg{25}{images/Tugas 3-214.png}
\begin{eulerformula}
\[
w=z^3-z
\]
\end{eulerformula}
\begin{eulerprompt}
>w=z^3-z;
>plot2d(w):
\end{eulerprompt}
\eulerimg{25}{images/Tugas 3-216.png}

\eulersubheading{Grafik Daerah Terbatas}
\begin{eulerformula}
\[
y=x^2
\]
\end{eulerformula}
\begin{eulerprompt}
>plot2d("x^2",-1,1,>filled,fillcolor=lightblue):
\end{eulerprompt}
\eulerimg{25}{images/Tugas 3-218.png}
\begin{eulerformula}
\[
y=x^2, \quad y=x^4
\]
\end{eulerformula}
\begin{eulerprompt}
>x=linspace(-1,1,100);
>plot2d(x,[x^2;x^4],>filled,fillcolor=cyan):
\end{eulerprompt}
\eulerimg{25}{images/Tugas 3-220.png}

\eulersubheading{Grafik Segi Banyak}
\begin{eulercomment}
Segitiga
\end{eulercomment}
\begin{eulerprompt}
>x=[0,2,1,0]; y=[0,0,2,0]; plot2d(x,y):
\end{eulerprompt}
\eulerimg{25}{images/Tugas 3-221.png}
\begin{eulerprompt}
>t=linspace(0,2*pi,7); x=cos(t); y=sin(t); plot2d(x,y,>filled,fillcolor=yellow):
\end{eulerprompt}
\eulerimg{25}{images/Tugas 3-222.png}

\eulerheading{Rujukan Lengkap Fungsi plot2d()}
\begin{eulerttcomment}
function plot2d (xv, yv, btest, a, b, c, d, xmin, xmax, r, n, ..
logplot, grid, frame, framecolor, square, color, thickness, style, ..
auto, add, user, delta, points, addpoints, pointstyle, bar, histogram, ..
distribution, even, steps, own, adaptive, hue, level, contour, ..
nc, filled, fillcolor, outline, title, xl, yl, maps, contourcolor, ..
contourwidth, ticks, margin, clipping, cx, cy, insimg, spectral, ..
cgrid, vertical, smaller, dl, niveau, levels)
\end{eulerttcomment}
\begin{eulercomment}
Fungsi plot serbaguna untuk plot pada bidang (plot 2D). Fungsi ini bisa melakukan plot untuk fungsi satu variabel, plot data, kurva pada bidang, plot batang, grid bilangan kompleks, dan plot implisit fungsi dua variabel.

\textbf{Parameter}

\begin{itemize}
    \item \texttt{x,y}: persamaan, fungsi, atau vektor data.
    \item \texttt{a,b,c,d}: area plot (bawaan \texttt{a=-2,b=2}).
    \item \texttt{r}: jika \texttt{r} diatur, maka \texttt{a=cx-r}, \texttt{b=cx+r}, \texttt{c=cy-r}, \texttt{d=cy+r}. \texttt{r} bisa berupa vektor \texttt{[rx,ry]} atau vektor \texttt{[rx1,rx2,ry1,ry2]}.
    \item \texttt{xmin,xmax}: rentang parameter untuk kurva.
    \item \texttt{auto}: menentukan rentang-$y$ secara otomatis (bawaan).
    \item \texttt{square}: jika \texttt{true}, mencoba mempertahankan rentang $x$-$y$ berbentuk persegi.
    \item \texttt{n}: jumlah interval (bawaan adalah adaptif).
    \item \texttt{grid}:
    \texttt{0} = tanpa grid dan label,
    \texttt{1} = sumbu saja,
    \texttt{2} = grid normal (lihat di bawah untuk jumlah garis grid),
    \texttt{3} = sumbu di dalam,
    \texttt{4} = tanpa grid,
    \texttt{5} = grid penuh termasuk margin,
    \texttt{6} = tanda sumbu pada bingkai,
    \texttt{7} = sumbu saja,
    \texttt{8} = sumbu saja, sub-tanda sumbu.
    \item \texttt{frame}: \texttt{0} = tanpa bingkai.
    \item \texttt{framecolor}: warna bingkai dan grid.
    \item \texttt{margin}: angka antara $0$ dan $0{,}4$ untuk margin di sekitar plot.
    \item \texttt{color}: warna kurva. Jika berupa vektor warna, warna tersebut akan digunakan untuk setiap baris matriks dari plot. Dalam kasus plot titik, parameter ini harus berupa vektor kolom. Jika vektor baris atau matriks warna penuh digunakan untuk plot titik, itu akan digunakan untuk setiap titik data.
    \item \texttt{thickness}: ketebalan garis pada kurva. Nilai ini bisa lebih kecil dari $1$ untuk garis yang sangat tipis.
    \item \texttt{style}: gaya plot untuk garis, penanda, dan isian.
    Untuk titik-titik gunakan \texttt{"[]"}, \texttt{"<>"}, \texttt{"."}, \texttt{".."}, \texttt{"..."}, \texttt{"*"}, \texttt{"+"}, \texttt{"|"}, \texttt{"-"}, \texttt{"o"}, \texttt{"[]\#"}, \texttt{"<>\#"}, \texttt{"o\#"} (bentuk terisi), \texttt{"[]w"}, \texttt{"<>w"}, \texttt{"ow"} (tidak transparan).
    Untuk garis gunakan \texttt{"-"}, \texttt{"--"}, \texttt{"-."}, \texttt{"."}, \texttt{".-."}, \texttt{"-.-"}, \texttt{"->"}.
    Untuk poligon terisi atau plot batang gunakan \texttt{"\#"}, \texttt{"\#O"}, \texttt{"O"}, \texttt{"/"}, \texttt{"\textbackslash{}"}, \texttt{"\textbackslash{}/"}, \texttt{"+"}, \texttt{"|"}, \texttt{"-"}, \texttt{"t"}.
    \item \texttt{points}: menggambar titik tunggal dan bukannya segmen garis.
    \item \texttt{addpoints}: jika \texttt{true}, menggambar segmen garis dan titik.
    \item \texttt{add}: tambahkan plot ini pada plot yang sudah ada.
    \item \texttt{user}: mengaktifkan interaksi pengguna untuk fungsi-fungsi.
    \item \texttt{delta}: ukuran langkah pada interaksi pengguna.
    \item \texttt{bar}: plot batang ($x$ adalah batas interval, $y$ adalah nilai interval).
    \item \texttt{histogram}: memplot frekuensi $x$ dalam $n$ subinterval.
    \item \texttt{distribution=n}: memplot distribusi $x$ dengan $n$ subinterval.
    \item \texttt{even}: menggunakan nilai bilangan bulat untuk histogram otomatis.
    \item \texttt{steps}: memplot fungsi sebagai fungsi tangga (\texttt{steps=1,2}).
    \item \texttt{adaptive}: gunakan plot adaptif ($n$ adalah jumlah langkah minimal).
    \item \texttt{level}: memplot garis level pada fungsi implisit dua variabel.
    \item \texttt{outline}: menggambar batas pada rentang level. Jika nilai levelnya berupa matriks $2 \times n$, rentang levelnya akan digambar dalam warna yang menggunakan gaya isian yang diberikan. Jika \texttt{outline} bernilai \texttt{true}, ia akan digambar dengan warna kontur. Dengan menggunakan fitur ini, area $f(x,y)$ di antara batas-batas dapat ditandai.
    \item \texttt{hue}: tambahkan warna \emph{hue} ke plot level untuk menunjukkan nilai fungsi.
    \item \texttt{contour}: gunakan plot level dengan level otomatis.
    \item \texttt{nc}: jumlah garis level otomatis.
    \item \texttt{title}: judul plot (bawaan \texttt{""}).
    \item \texttt{xl, yl}: label untuk sumbu-$x$ dan sumbu-$y$.
    \item \texttt{smaller}: jika $>0$, akan ada lebih banyak ruang ke kiri untuk label.
    \item \texttt{vertical}: mengaktifkan atau menonaktifkan label vertikal. Ini mengubah variabel global \texttt{verticallabels} secara lokal untuk satu plot. Nilai \texttt{1} mengatur teks hanya vertikal, nilai \texttt{2} menggunakan label numerik vertikal pada sumbu-$y$.
    \item \texttt{filled}: mengisi plot pada sebuah kurva.
    \item \texttt{fillcolor}: warna isian untuk plot batang dan kurva terisi.
    \item \texttt{outline}: batas untuk poligon terisi.
    \item \texttt{logplot}: mengatur plot logaritmik (\texttt{1} = logplot di $y$, \texttt{2} = logplot di $xy$, \texttt{3} = logplot di $x$).
    \item \texttt{own}: sebuah string yang menunjuk pada rutin plot sendiri. Dengan \texttt{>user}, Anda mendapatkan interaksi pengguna yang sama seperti pada \texttt{plot2d}. Rentangnya akan ditetapkan sebelum tiap panggilan ke fungsi Anda.
    \item \texttt{maps}: memetakan ekspresi (\texttt{0} lebih cepat), fungsi selalu dipetakan.
    \item \texttt{contourcolor}: warna garis kontur.
    \item \texttt{contourwidth}: lebar garis kontur.
    \item \texttt{clipping}: beralih ke \emph{clipping} (bawaannya adalah \texttt{true}).
    \item \texttt{cgrid}: menentukan jumlah garis grid untuk plot grid kompleks. Harus merupakan pembagi dari ukuran matriks dikurangi $1$ (jumlah subinterval). \texttt{cgrid} bisa berupa vektor \texttt{[cx,cy]}.
\end{itemize}

\textbf{Gambaran Umum}

Fungsi ini dapat memplot:
\begin{itemize}
    \item ekspresi, kumpulan panggilan (\emph{call collections}), atau fungsi dari satu variabel,
    \item kurva parametrik,
    \item data $x$ terhadap data $y$,
    \item fungsi implisit,
    \item plot batang,
    \item grid kompleks,
    \item poligon.
\end{itemize}

Jika fungsi atau ekspresi untuk \texttt{xv} diberikan, \texttt{plot2d()} akan menghitung nilai dalam rentang yang diberikan menggunakan fungsi atau ekspresi tersebut. Ekspresi harus berupa ekspresi dalam variabel $x$. Rentang harus didefinisikan pada parameter \texttt{a} dan \texttt{b} kecuali rentang bawaan $[-2,2]$ yang akan digunakan. Rentang-$y$ akan dihitung secara otomatis, kecuali \texttt{c} dan \texttt{d} ditentukan, atau suatu radius \texttt{r}, yang akan menghasilkan rentang $[-r,r]$ untuk $x$ dan $y$. Untuk plot dari fungsi-fungsi, \texttt{plot2d} akan menggunakan evaluasi adaptif pada fungsi tersebut secara bawaan. Untuk mempercepat plot pada fungsi yang rumit, matikan hal ini dengan \texttt{<adaptive}, dan secara opsional turunkan jumlah interval \texttt{n}. Selain itu, \texttt{plot2d()} secara bawaan akan menggunakan pemetaan. Maksudnya, ia akan menghitung elemen plot demi elemen. Jika ekspresi Anda atau fungsi Anda bisa menangani vektor $x$, Anda dapat mematikannya dengan \texttt{<maps} agar evaluasi berjalan lebih cepat.

Perhatikan bahwa plot adaptif selalu dihitung elemen demi elemen. Jika fungsi atau ekspresi untuk keduanya baik \texttt{xv} maupun \texttt{yv} ditentukan, \texttt{plot2d()} akan menghitung kurva dengan nilai-nilai \texttt{xv} sebagai koordinat-$x$ dan nilai-nilai \texttt{yv} sebagai koordinat-$y$. Dalam kasus ini, rentang harus didefinisikan untuk parameter tersebut menggunakan \texttt{xmin}, \texttt{xmax}. Ekspresi yang terkandung di dalam string harus selalu berupa ekspresi dalam variabel parameter $x$.
\end{eulercomment}

\end{eulernotebook}
\end{document}